% Sablon pentru realizarea lucrarii de licenta, conform cu recomandarile
% din ghidul de redactare:
% - https://fmi.unibuc.ro/finalizare-studii/
% - https://drive.google.com/file/d/1xj9kZZgTkcKMJkMLRuoYRgLQ1O8CX0mv/view

% Multumiri lui Gabriel Majeri, acest sablon a fost creat pe baza
% codului sursa a lucrarii sale de licenta. 
% Codul sursa: https://github.com/GabrielMajeri/bachelors-thesis
% Website: https://www.gabrielmajeri.ro/
%
% Aceast sablon este licentiat sub Creative Commons Attribution 4.0 International License.

\documentclass[12pt, a4paper]{report}

% Suport pentru diacritice și alte simboluri
\usepackage{fontspec}

% Suport pentru mai multe limbi
\usepackage{polyglossia}

% Setează limba textului la română
\setdefaultlanguage{romanian}
% Am nevoie de engleză pentru rezumat
\setotherlanguages{english}

% Indentează și primul paragraf al fiecărei noi secțiuni
\SetLanguageKeys{romanian}{indentfirst=true}

% Suport pentru diferite stiluri de ghilimele
\usepackage{csquotes}

\DeclareQuoteStyle{romanian}
  {\quotedblbase}
  {\textquotedblright}
  {\guillemotleft}
  {\guillemotright}

% Utilizează biblatex pentru referințe bibliografice
\usepackage[
    maxbibnames=50,
    sorting=nty
]{biblatex}

\addbibresource{bibliography.bib}

% Setează spațiere inter-linie la 1.5
\usepackage{setspace}
\onehalfspacing

% Modificarea geometriei paginii
\usepackage{geometry}

% Include funcțiile de grafică
\usepackage{graphicx}
% Încarcă imaginile din directorul `images`
\graphicspath{{./images/}}

% Listări de cod
\usepackage{listings}

% Linkuri interactive în PDF
\usepackage[
    colorlinks,
    linkcolor={black},
    menucolor={black},
    citecolor={black},
    urlcolor={blue}
]{hyperref}

% Simboluri matematice codificate Unicode
\usepackage[warnings-off={mathtools-colon,mathtools-overbracket}]{unicode-math}

% Comenzi matematice
\usepackage{amsmath}
\usepackage{mathtools}

% Formule matematice
\newcommand{\bigO}[1]{\symcal{O}\left(#1\right)}
\DeclarePairedDelimiter\abs{\lvert}{\rvert}

% Suport pentru rezumat în două limbi
% Bazat pe https://tex.stackexchange.com/a/70818
\newenvironment{abstractpage}
  {\cleardoublepage\vspace*{\fill}\thispagestyle{empty}}
  {\vfill\cleardoublepage}
\renewenvironment{abstract}[1]
  {\bigskip\selectlanguage{#1}%
   \begin{center}\bfseries\abstractname\end{center}}
  {\par\bigskip}

% Suport pentru anexe
\usepackage{appendix}

% Stiluri diferite de headere și footere
\usepackage{fancyhdr}

% Metadate
\title{Implementarea arhitecturii bazate pe microservicii în Kubernetes}
\author{Bahrim Dragoș}

% Generează variabilele cu @
\makeatletter

\begin{document}

% Front matter
\cleardoublepage
\let\ps@plain

% Pagina de titlu
\begin{titlepage}

% Redu marginile
\newgeometry{left=2cm,right=2cm,bottom=1cm}

\begin{figure}[!htb]
    \centering
    \begin{minipage}{0.2\textwidth}
        \includegraphics[width=\linewidth]{logo-ub.png}
    \end{minipage}
    \begin{minipage}{0.5\textwidth}
        \large
        \vspace{0.2cm}
        \begin{center}
            \textbf{UNIVERSITATEA DIN BUCUREȘTI}
        \end{center}
        \vspace{0.3cm}
        \begin{center}
            \textbf{
                FACULTATEA DE \\
                MATEMATICĂ ȘI INFORMATICĂ
            }
        \end{center}
    \end{minipage}
    \begin{minipage}{0.2\textwidth}
        \includegraphics[width=\linewidth]{logo-fmi.png}
    \end{minipage}
\end{figure}

\begin{center}
\textbf{SPECIALIZAREA TEHNOLOGIA INFORMAȚIEI}
\end{center}

\vspace{1cm}

\begin{center}
\Large \textbf{Proiect de diplomă}
\end{center}

\begin{center}
\huge \textbf{\MakeUppercase{\@title}}
\end{center}

\vspace{3cm}

\begin{center}
\large \textbf{Absolvent \\ \@author}
\end{center}

\vspace{0.25cm}

\begin{center}
\large \textbf{Coordonator științific \\ Conferențiar Doctor Kevorchian Cristian}
\end{center}

\vspace{2cm}

\begin{center}
\Large \textbf{București, iunie-iulie 2023}
\end{center}
\end{titlepage}
\restoregeometry
\newgeometry{
  margin=2.5cm
}

\fancypagestyle{main}{
  \fancyhf{}
  \renewcommand\headrulewidth{0pt}
  \fancyhead[C]{}
  \fancyfoot[C]{\thepage}
}

\addtocounter{page}{1}

% Rezumatul
\begin{abstractpage}

\begin{abstract}{romanian}
Microserviciile sunt o arhitectură de sisteme distribuite ce a luat amploare în urma 
structurării organizațiilor către piață. 
Aceasta a fost facilitată de schimbări în modul în care lansăm produsele prin 
dezvoltarea platformelor cloud și a ce a dus la îmbunătățiri la timpul de 
livrare dar și a consistenței. Obiectivul este prezentarea acestor concepte și 
implementărea acestora prezentând modul de gândire ce influențează ciclul produsului. 
În acest scop, mă documentez legat de domeniu și încerc 
să îmi fac o idee de ansamblu asupra aspectelor ce trebuie luate în construirea unui sistem 
iar apoi o să aplic aceste informații. Aceasta lucrare ar trebui să servească ca un prim 
contact cu microserviciile dar și a practicilor DevOps.
\end{abstract}

\begin{abstract}{english}
Lorem ipsum dolor sit amet, consectetur adipiscing elit. Fusce vitae eros sit amet sem ornare varius. Duis eget felis eget risus posuere luctus. Integer odio metus, eleifend at nunc vitae, rutrum fermentum leo. Quisque rutrum vitae risus nec porta. Nunc eu orci euismod, ornare risus at, accumsan augue. Ut tincidunt pharetra convallis. Maecenas ut pretium ex. Morbi tellus dui, viverra quis augue at, tincidunt hendrerit orci. Lorem ipsum dolor sit amet, consectetur adipiscing elit. Aliquam quis sollicitudin nunc. Sed sollicitudin purus dapibus mi fringilla, nec tincidunt nunc eleifend. Nam ut molestie erat. Integer eros dolor, viverra quis massa at, auctor.
\end{abstract}

\end{abstractpage}

\tableofcontents

% Main matter
\cleardoublepage
\pagestyle{main}
\let\ps@plain\ps@main

\chapter{Introducere}
Domeniul informaticii, calculatoarelor și al tehnologiei informației este unul care va fi mereu
în continuă evoluție. Dorința companiilor de a-și servi clienții cât mai rapid a dus la un avans
tehnologic în care mereu apare ceva nou ce are ca scop îmbunătățirea proceselor actuale. Această arie și-a
început dezvoltarea aproape acum o sută de ani iar avansul poate fi apreciat numai gândidu-ne la modul în
care scriam cod și livram produse acum, acum zece ani și acum două zeci de ani. Se observă o diferență
și la accesibilitatea unităților de calcul, actual majoritatea persoanelor au acces la un dispozitiv ce
se poate conecta la Internet, ceea ce le permite să devină consumatori la diferite servicii.

Piața pentru unități de calcul a evoluat. În 1943, Thomas Watson, director la IBM menționează că 
„I think there is a world market for maybe five computers” (Cred că există o piață globală 
pentru poate cinci computere). În 1999, Michael Barr afirmă în „Programming Embedded Systems in C and C++”
„One of the more surprising developments of the last few decades has been the
ascendance of computers to a position of prevalence in human affairs. Today there are
more computers in our homes and offices than there are people who live and work in
them. Yet many of these computers are not recognized as such by their users.” (Una dintre cele mai 
surprinzătoare dezvoltări a ultimelor decenii a fost ascensiunea computerelor într-o posiție predominantă 
în afacerile oamenilor. Astăzi sunt mai multe computere în casele noastre și în birouri decât persoane care trăiesc
și lucrează în ele. Însă multe dintre acestea nu sunt recunoscute de către utilizatorii lor). Cele două fraze evidențiază
modul cum s-a schimbat prezența computerelor în viața noastră. La început acestea erau foarte costisitoare, 
greu de administrat și de operat însă acestea au devenit din ce în ce mai mici iar recent aproape orice are 
nevoie de semiconductori datorită integrării ce le fac „smart”. 

Aceste afirmații se aplică și pentru produsele software, ce au ca menire să ofere asistență altor aplicații
sau să fie consumate direct. Indiferent de locul în care mă duc, probabil persoana cu care aș interactiona
folosește un computer, fie că îmi cumpăr ceva de la un magazin făcând o plată cu cardul sau că mi se livrează un colet
iar curierul marchează pe AWB că a fost livrat. Christopher Little afirmă că „Every company is a technology company, regardless of 
what business they think they are in. A bank is just an IT company with a banking license.” (Orice companie
este o companie ce se axează pe tehnologie, indiferent de mediul de afaceri în care se află. O bancă este doar o companie
IT cu licență de a funcționa ca o bancă.)

Desigur nu toate companiile sunt la fel și nu toate concurează cu toate companiile, însă chiar dacă 
nu ești o companie ce activează în tehnologie, probabil folosești componente tehnologice pentru a-ți
îmbunătății randamentul sau să oferi clienților un avantaj față de competiție, care probabil gândește în
același mod.

Însă pentru fiecare companie obiectivul este identic, să livreze clienților produse de calitate și cât mai 
rapid. Pentru aceasta trebuie să ne asigurăm că avem inginerii necesari pentru implementarea cererilor noi,
însă momentul în care aceasta este terminată este doar începutul procesului de integrare și ulterior lansare,
și uneori acesta este cel care reprezintă un blocaj din mai multe puncte de vedere ce limiteaza numarul de câte
ori putem aduce o îmbunătățire produsului nostru.

În încercarea de a livra mai rapid sau pentru a ușura modul în care manevrăm un influx sau o lipsă de trafic,
s-a încercat folosirea unei arhitecturi paralele asupra aplicațiilor creând arhitecturi orientate pe servicii, 
iar în momentul în care aceste servicii sunt concentrare pe un număr redus de funcționalități și pot fi lansate
independent și să își păstreze funcționalitate putem vorba de o arhitectură bazata pe microservicii.

Însă schimbând arhitectura sau practiciile din modul de desfășurare al livrarii produsului nu este suficient
pentru a asigura performanța, întrucât ambele dintre ele aduc dezavantaje ce trebuie tratate iar uneori 
acestea sunt mai mari decât avantajele pe care schimbarea le-ar aduce, însă implementarea corectă în locurile
ce ar beneficia de o astfel de abordare poate să aducă performanțe ce nu ar fi posibile cu procesele vechi.

Afinitatea mea pentru microservicii provine de la lipsa de cunoștiințelor pentru a putea alege dintre mai
multe limbaje de programare, framework-uri sau platforme. De exemplu, ce avantaje mi-ar face să aleg ca 
pentru server-ul meu să folosesc o aplicație în Node.JS sau una in Go? Chiar dacă putem să cunoaștem 
niște lucruri informative la început, cei mai buni indicatori sunt monitorizările proprii în producție. De 
asemenea pot fi cazuri în care îmbunătățirile aduse unei platforme pe parcurs o fac mai bună față de ce era
inițial (dezvoltarea TypeScript poate să influențeze alegerea inițială a unei platforme SpringBoot ce este type-safe).
Astfel, pe o arhitectură bazată pe microservicii putem să folosim limbajul potrivit pentru serviciul căruia
îi aduce cele mai multe beneficii. Acest lucru se extinde și pentru baze de date. Poate în timpul evoluției 
serviciului apar tehnologii noi, însă având un singur lucru central ar face migrarea mult mai grea, de asemenea
uneori ar trebui să facem compromisuri pe care alegerea unui alt tip de baze de date ar face să o dispară.

Lucrarea mea va urma o structură în care aprofundez aspectele teoretice la început iar la final încerc să
aplic elementele prezentate în crearea unei aplicații. Înițial pornesc cu prezentarea microserviciilor
continuând cu diferitele moduri în care acestea pot fi lansate, iar
la final o prezentare generală a modului în care îmi construiesc aplicația.

\chapter{Preliminarii}
Cel puțin până la momentul la care am decis ce temă să aleg pentru lucrarea mea de licență niciun curs de bază 
nu a abordat în detaliu tema pe care vreau să o dezvolt, din acest motiv lucrarea mea de licență nu este doar
aprofundarea unui concept și crearea unei aplicații, ci încercarea familiarizării cu obiectele de lucru în același timp.
Desigur, unele aspecte s-ar putea desprinde și în urma practicii în industrie dar nu a fost suficient ca să
pot să consider că am cunoștințe de bază în acest domeniu și doar încerc să găsesc lucruri cât mai specifice.

De asemenea, microserviciile sunt o arhitectură destul de nouă ce a luat amploare datorită limitărilor din
punctul de vedere al dezvoltării și scalării aplicațiilor monolitice, este destul
greu să vizualizezi anumite lucruri fără experiență. Întrucât produsul
pe care îl dezvolt nu ar necesita neapărat o arhitectură bazată pe microservicii, ar funcția complet normal
și aș reduce foarte mult din complexitate dacă ar fi dezvoltată în mod normal.
\chapter{Microservicii}
\section{Concepte generale}
Microserviciile sunt o arhitectură de sisteme distribuite ce a devenit populară în ultimele decenii datorită
necesității scalării. Ceea ce o separă de o arhitectură orientate pe servicii este faptul că microserviciile
sunt mult mai specifice, acoperind o singură parte din afacere și faptul că pot fi lansate în execuție independent,
astfel microserviciile ar trebui să depindă cât mai puțin de altele. Comunicarea se face direct către interfețele
expuse de fiecare, iar informațiile reținute de fiecare (de exemplu baza de date de date sau diferite fișiere)
pot să fie modificate doar prin metodele de comunicare oferite ci nu direct.

Un concept important în microservicii este ascunderea de informații, astfel din exterior fiecare ar trebui să 
fie tratat ca o cutie neagră ce face anumite lucruri și expune anumite date, nu contează detaliile de implementare
și nici tehnologiile folosite în implementare, ci doar capabilitățile fiecăruia. Acest lucru permite să modificăm
microserviciul în funcție de cerințe, chiar și refactorizarea metodelor inițiale cât timp se respectă 
capabilitățile pe care microserviciul ar trebui să le îndeplinească.

Cel mai important aspect al microserviciilor este capacitatea de a lucra independent de celelalte. Acestea
trebuie să funcționeze independent de celelalte iar dacă aducem schimbări într-unul atunci nu suntem obligați să
facem schimbări în alt microserviciu, în caz contrat ar apărea dificultăți în lansare. Obținerea independenței poate
să aducă contribuții majore în timpul de deployment însă implementarea este mult mai complicată datorită necesității
comunicării cu alte servicii.

Un alt aspect de bază este acoperirea unui singur segment din afacere prin intermediul unui microserviciu. Alături 
de caracteristica precedentă, dacă afacerea și-ar dezvolta necesitățile ar fi mult mai dificil să modificăm
mai multe microservicii și ulterior să coordonăm lansarea lor.

Independența microserviciilor este importantă, în acest scop, ele ar trebui să fie singurele care ar putea să își modifice starea.
Întrucât starea se referă la elementele componente, de exemplu o bază de dată, acestea nu ar trebui să fie împărtășite
și nici accesate de alte microservicii. De asemenea ar trebui evitate modificarea interfețelor declarate pentru
a evita necesitatea modificării altor microservicii.
„
Dimensiunea microserviciile ar trebui să fie redusă prin oferirea unui număr restrâns de funcționalități declarate pentru a
a fi ușor de înțeles și întreținut. Însă important este cât de multe microservicii
pot fi administrate întrucât un număr mai mare de microservicii, deși mici ca funcționalități, aduc dificultăți
în comunicare, scalare și depanare.

Microserviciile sunt agnostice din punct de vedere al tehnologiilor care ar putea fi folosite. Din acest motiv, acestea
oferă flexibilitate însă prețul plătit este crearea a mai multor puncte vulnerabile.

Orice prezentare al acestui tip de arhitectură nu ar fi completă fără prezentarea conceptului de „monolith” (monolit), întrucât
este considerată o arhitectură veche reprezentând modul de funcționare din trecut. În cel mai simplist mod de funcționare, 
sistemul monolitic este reprezentat de un singur proces „gigantic” ce acoperă toată funcționalitatea sistemului, în general conectat
la o singură bază de date ce și ea este reprezentată de un singur proces. Putem intui problemele cu această organizare,
dacă conexiunea către procesul monolitic cade sau baza de date are probleme, întreg sistemul este afectat. Microserviciile incearcă 
să acopere aceste probleme prin independență, astfel dacă un microserviciu este căzut, celelalte nu sunt afectate.
Procesul poate fi însă separat pe module dezvoltate separat însă care la final se cuplează formând un singur proces.
Există conceptul de sistem monolitic distribuit în care acesta este la fel separat în module independente însă care nu ar funcționa
dacă nu sunt toate active, acest tip de sistem oferă toate dezavantajele sistemelor distribuite dar și a monolitului ceea 
ce îl face destul de rar întâlnit.
                  
Chiar dacă o arhitectură monolitică prezintă dezavantaje evidente ce sunt ușor de văzut, acestea nu sunt la fel de mari
pentru companiile mici. Numărul redus de aplicații ce trebuie lansate și monitorizate,
posibilitatea de reutilizare a codului mult mai ușoară, obținerea infrastructurii pentru
proces și uneori chiar și scalarea, deși este limitată poate să facă o arhitectură monolitică
mult mai atractivă. Nu ar trebui să asociem arhitectura monolitică ca un lucru antic, ci să 
o considerăm ca o opțiune. Însă pentru o organizație mică, acesta ar trebui să fie doar un
punct de plecare, eventual după ce serviciul oferit devine mai popular s-ar putea să ajungem
să cunoaștem limitările arhitecturii și să ne orientăm tot către microservicii, însă dacă 
nu avem succes, eliminăm multe complexități apărute în urma introducerii microserviciilor.

Principalele avantaje ale microserviciilor provin de la baza acestora, o arhitectură distribuită,
însă cuplat cu conceptul de ascundere a informației și domain-driven design (design orientat domeniu)
aduc multe alte avantaje asupra altor arhitecturi distribuite. În cadrul unui sistem format din microservicii
putem folosi orice tip de tehnologie pentru program și orice tip de bază de date, întrucât
ascundem implementarea de exterior. Astfel putem alege tehnologii ce aduc avantaje în 
dezvoltarea microserviciului. Acest avantaj duce la capabilitatea de a folosi tehnologii noi
fară a afecta tot sistemul, însă limitează capacitatea de a împărtășii cod (de exemplu prin
librării interne, întrucât acestea ar trebui să fie rescrice pentru fiecare limbaj). 
Un sistem distribuit rezistă mult mai ușor la căderi, întrucât acestea pot fi tratate ca să
se prevină un lanț, însă cu un număr mare de puncte slabe devine greu de aproximat cât de
mult este afectat un sistem. Scalarea unui sistem monolitic este ușoară, doar replicăm procesul,
însă la microservicii scalarea poate fi mult mai concentrată pe punctele care chiar au nevoie să fie
scalate. Lansarea unei schimbări într-un microserviciu este mult mai ușoară întrucât nu necesită
crearea de timp mort în aplicație, întrucât nu tot sistemul este înlocuit ci doar o mică parte din acesta,
ce poate fi chiar și mai controlată prin diferite aplicații de orchestrare.

Adaptarea microserviciilor vine cu un cost, pe lângă dificultățile de implementare a design-ului apar
și alte probleme, de exemplu dezvoltarea locală. Atunci când dezvoltăm un microserviciu acesta, posibil,
este nevoit să comunice cu alte microservicii care la rândul lor comunică cu alte microservicii și putem continua,
toate acestea microservicii trebuie să funcționeze în același timp pe laptop-ul dezvoltatorului,
întrucât nu putem folosi microserviciile din producție, însă nu e posibil să rulăm foarte multe
microservicii pe un singur computer depinzând de configurație. Adoptarea microserviciilor,
în general nu poate avea succes decât dacă este combinată cu elemente de DevOps, ceea ce înseamnă
ca dezvoltatorii și operatorii să învețe tehnologii noi, ceea ce este un împediment în momentul
în care vrem să lansăm un produs rapid în piață și nu avem experiența necesară. Microserviciile necesită
mult mai multă infrastructură pentru a fi rulată, mai multă tehnologie ce trebuie să fie folosită
pentru ca dezvoltarea și livrarea să decurgă eficient, însă pe termen lung acesta poate să
fie prevăzută dacă arhitectura este folosită cum trebuie prin livrare mai rapidă și mai eficientă
deci profitul ar putea să crească. Monitorizarea sistemelui devine mai complicată, într-un proces monolitic
toate fișierele de jurnalizare (logging) sunt în același loc. Într-un sistem distribuit, acestea 
trebuie să fie colectate, și ulterior ansamblate pentru a avea sens, acest lucru devine mai dificil când
intră în discuție fenomenul de replicare al unei instanțe. Securizarea unui sistem distribuit este
mai grea, traficul de date se face mult mai mult prin rețele sau prin diferite căi alternative ceea 
ce necesită un grad ridicat de atenție. Testarea unui sistem format din microservicii este mult mai
dificilă, mai ales când trebuie testate mai multe microservicii în același timp, acesta duce la creșterea
timpului necesar pentru a primi răspuns. Sistemele distribuite cresc timpul de răspuns întrucât acestea nu
se mai face în cadrul unui singur proces, deci datele trebuie codificate și decodificate atunci când trec
dintr-un mediu în altul. În cadrul unui sistem distribuit o altă problemă care apare este consistența 
datelor, dată de faptul că în timpul unei cereri, datele dintr-un microserviciu se pot schimba, astfel
tranzacțiile devin dificile mai ales când înainte ne bazam pe atomicitatea și consistența oferită de 
baza de date.

Astfel, microserviciile funcționează cel mai bine pentru organizațiile mari ce vor ca dezvoltatorii
lor să livreze facilități noi foarte repede care să nu fie întârziate de compatibilitatea cu 
caracteristici deja existente. Prin natura variată a acestei arhitecturi se pot folosi tehnologii noi
dar și folosirea platformelor cloud, astfel se oferă multă flexibilitate cu prețul creșterii 
complexității. Pentru companiile mici, al căror rată de succes este greu de calculat, overhead-ul 
produs de adaptarea unei arhitecturi în acest stil nu ar aduce suficient de multe avantaje și ar
trebui să opteze pentru o arhitectură tradițională ce este stabilă și oferă colaborare mult mai ușoară
în cadrul aceluiași proiect la scale mici. 

\section{Design}  
Partea de creare a schemei arhitecturale a unui sistem este una dintre
cele mai importante faze ale dezvoltării. Acesta este momentul în care putem descoperii
că microserviciile nu ar fi chiar o alegere potrivită pentru cerințele noastre. Din acest
motiv este important să nu ignorăm potențiale probleme ci să se formuleze toate problemele
pe care le putem aștepta în implementare

\subsection{Interconectare}
Deși microserviciile pot fi lansate și pot funcționa independent, acestea trebuie să funcționeze
în cadrul unui sistem complex. Astfel trebuie definite modul de comunicare și interfețele expuse.
În definirea interconectării întâlnim câteva concepte de bază prezentate anterior ce provin
din descompunerea în module.

Ascunderea de informații este procesul de separare al detaliilor de implementare pentru a evita
modificarea în alte locuri în cazul în care acestea se schimbă. Acest concept ofer varii beneficii
precum îmbunătățirea timpului de dezvoltare datorită posibilități de dezvoltare în paralel,
creșterea întelegerii sistemului mult mai ușor întrucât fiecare modul poate să fie înțeles
independent și flexibilitate oferit de faptul că putem schimba funcționalitate fără a fi necesar
ca alte părți ale sistemului să fie modificate. David Parnas spune că „The connections between modules are the assumptions which the modules make about
each other.” (Conexiunile create între module sunt ipotezele create de module față de celelalte) \cite{DBLP:conf/ifip/Parnas71}.
Astfel, cu cât alte microservicii folosesc un anumit microserviciu mai mult se crează mai multe
locuri în care devine mai dificil de modificat, astfel cu cât întâmpinăm problema aceasta mai puțin
cu atât putem lansa acel microserviciu fără a lua în considerare modificări care să pastreze compatibilitatea
cu elementele existente.

Un alt concept important este coeziunea, scopul microserviciilor este de a funcționa independent,
astfel vrem ca lucrurile care înteracționează să fie grupate iar atunci când se schimbă ceva 
să nu trebuiască să modificăm multe locuri. Elementele care nu interacționează ar trebui să fie separate
pentru a nu crea probleme de regresie.

În sisteme distribuite apare cuplarea, atunci când un microserviciu depinde mai mult sau mai puțin de
altul pentru a avea funcționalitate completă. Un sistem cu cuplare mică poate să fie lansat fără să trebuiască
modificări în locuri multiple.

Cele două concepte definesc lucruri asemănătoare, astfel ele se pot rezuma ușor ca „A structure is stable if cohesion is strong and coupling is low.” \cite{endres_rombach_2003} 
(O structură este stabilă dacă coeziunea este ridicată și cuplarea este scăzută). Coeziunea se referă la lucrurile
din interior iar cuplarea se referă la exterior iar acest lucru este greu de balansat.

Cuplarea poate apărea sub diferite forme, cuplare de domeniu în care un microserviciu interacționează
cu alt microserviciu pentru că are nevoie de funcționalitate ce depăsește atribuțiile lui. Astfel de 
cuplare, în general, nu poate fi evitată într-un sistem distribuit. Acesta este cel mai scăzut nivel
de cuplare însă trebuie să evităm cazul în care un microserviciu depinde un centru pentru alte microservicii
întrucât ar însemna ca acesta are prea multe responsabilități. Există conceptul de cuplare temporară, adică
atunci când un microserviciu așteaptă ca alt microserviciu să execute ceva, astfel cuplarea se întâmplă pentru că cele
două operațiuni se întâmplă în același timp. Cuplarea Pass-Through apare atunci când un microserviciu este nevoit
să trimită date unui alt microserviciu întrucât acesta trebuie să trimite datele 
unui microserviciu la un nivel mai mic, astfel un microserviciu are rolul doar de a transmite date. Acest lucru
este problematic întrucât dacă formatul acestor date ar trebui să se schimbe, ar trebui să
modificăm mai multe microservicii. Evitarea acestui tip de cuplare se face prin schimbarea modului
prin care comunicăm de exemplu microserviciul ce emite date ar putea să trimită direct datele către microserviciul ce are
nevoie de ele sau ca microserviciul interpediar să trimită datele in format pur, fără să îl intereseze o formă anume.
Cuplarea comună apare atunci când mai multe microservicii încearcă să acceseze o resursă 
comună, de exemplu o bază de date sau locația unui fișier. Problema în acest caz este că dacă modul în care
această resursă interacționează se schimbă, mai multe microservicii trebuie modificate. Sau dacă
unul din microservicii are nevoie de un comportament special, alte microservicii ar trebui să 
se adapteaze ceea ce crează probleme în dezvoltare. Soluții posibile de evitare pot fi crearea unui alt 
microserviciu ce funcționează ca o interfață pentru resursa comună ce va funcționa dupa principiile unui
microserviciu. Cuplarea de conținut apare atunci când un microserviciu modifică starea internă a altui 
microserviciu, de exemplu prin modificarea bazei de date a acesteia. Acest lucru este periculos întrucât pot
apărea modficări nedorite prin utilizare incorectă, astfel se încalcă procesul de ascundere de informații.

\subsection{Domain Driven Design}
Domain Driven Design este un concept introdus de Eric Evans în cartea sa \cite{ddd}. Acesta se bazează
pe crearea codului în jurul afacerilor pentru ca acesta să fie mai ușor de înțeles prin creșterea
interacțiunii cu analiștii de afaceri, acest lucru este benefic pentru microservicii întrucât acestea se 
pot orienta către mediul de afaceri pentru a se dezvolta alături de acesta fiind mult mai eficient.

Unul din concepte este limbajul universal ce cuprinde ca codul din interiorul unui microserviciu să 
folosească termeni folosiți de către utilizatori și de către analiști, acest lucru este benefic întrucât
face comunicarea mult mai eficientă, o persoană nou venită poate să înțeleagă afacerea și codul în același timp.
Dacă cele două sunt diferite, atunci pot exista probleme în comunicarea problemelor.

În DDD, un agregat este o unitate ce are o anumită stare ce se poate modifică în timp și are ca scop
modelarea unui concept din domeniul real. În general un agregat este cuprins în întregime de către
un microserviciu chiar dacă acesta poate conține mai multe microservicii. Un agregat poate fi schimbat
din exterior însă poate să își controleze singur starea pentru a preveni stările ilegale

Mărginirea contextului se referă la ascunderea mai multor activități ce țin de implementare
cu scopul de a defini anumite responsabilități sigure pe care le îndeplinește. Acest lucru se
poate îndeplini prin ascunderea modelului, atunci când vrem să expunem o entitate nu trebuie să
oferim toate informațiile despre acesta ci doar cele necesare. Distribuirea unui model, adică a
unei entități între mai multe microservicii se face prin traducerea acestei entități în domeniu propriu.

DDD în microservicii este un concept fundamental întrucât ambele pun la centru afacerea și
dezvoltarea în jurul ei. Astfel conceptele prezentate în DDD pot fi aplicate microserviciilor
mărginirea contextului are ca bază ascunderea de informații și duce la crearea de formate standard.

\subsection{Concepte de programare orientată obiect în microservicii}

Multe din problemele cu care ne confruntăm în microservicii pot fi tratate folosind concepte
ce au ca scop ușurarea modului în care creăm sisteme prin intermediul claselor în programarea 
orientată obiect. Deși diferite în natură, ambele au ca rol definirea unor interacțiuni între
elemente ce sunt făcute să funcționeze separat sau prin intermediul unuor interfețe.

Începând cu conceptele de bază, abstractizarea se referă la ascunderea implementărilor interne
și afișarea doar facilitățile oferite de clasă ce pot fi accesate din exterior. Acest lucru face 
referire la conceptul de ascundere al informațiilor descris anterior.

Encapsularea se referă la agregarea de date și atribute, adică funcționalități pentru a crea un tot
unitar, acest principiu este compatibil cu microserviciile și nevoia acestora de a funcționa
independent folosind datele proprii și permite o securizare mai bună a stării interne.

Moștenirea se referă la capacitatea unei clase de a prelua atributele și funcționalitățile unei clase
deja existente, deși mult mai greu de imaginat în microservicii, moștenirea în microservicii ne permite
să creăm microservicii ce extind pe cel existent fără a cauza probleme.

Polimorfismul este un concept ce face referire la capacitatea accesarii a unor multiple entități
prin intermediul unei singure interfețe, în contextul microserviciilor putem spune că un microserviciu
ar putea să înlocuiască alt microserviciu dacă are aceași interfață de comunicare și duce la
crearea a unor aceleași rezultate.

Pentru microservicii putem aplica și concepte mai complicate de OOP precum principiile SOLID.
Principiile SOLID sunt un set de reguli cu scopul de a oferi dezvoltatorilor o metodă de a
scrie cod, astfel încât acesta să respecte cerințele actuale și să permită dezvoltări/modificări
ulterioare cu ușurință fără a avea dificultăți cu metodele deja existente. SOLID este un acronim
iar fiecare literă face referire la un principiu. Principiul responsabilității unice (Single Responsible
Principle) se referă la necesitatea ca fiecare clasă și metodă să aibă o singură funcționalitate,
făcând asta reducem complexitatea dar și eventuale redundanțe create prin folosirea aceluiași
cod în mai multe locuri făcând ca fiecare porțiune din cod să aibă un singur motiv de a exista.
Principiul Deschis – Închis (Open Close Principle) conduce la faptul că entitățile scrise trebuie să
fie deschise pentru moștenire dar să nu poată fie modificate, astfel asemănător cu principiul
anterior, codul deja scris nu va fi afectat întrucât noile modificări se aduc asupra unei porțiuni
noi. Principiul substituției al lui Liskov (Liskov Substitution Principle) se referă la faptul că clasele
derivate trebuie să acopere funcționalitățile clasei pe care o moștenesc astfel încât acestea să
poată să le înlocuiască fără modificări adiționale. Principiul separării interfețelor (Interface
Segregation Principle), clasele ce implementează interfețe nu ar trebui să implementeze metode
de care nu au nevoie, astfel dacă este nevoie să adăugăm o metodă nouă unei interfețe aceasta
trebuie să fie implementată doar de către obiectele care au nevoie, dacă acest fapt nu este
satisfăcut trebuie ca interfețele să fie separate ca să nu forțăm alte clase să implementeze
metode noi. Principiul inversării dependințelor (Dependency Inversion Principle) se referă că
modulele de nivel înalt nu ar trebui să depindă de cele de nivel scăzut, ele trebuie să fie
abstractizate astfel acestea pot fi modificate în orice moment dacă vrem să aducem un alt modul
cu funcționalități similare dar cu implementări diferite.\cite{temasolid}

\subsection{Migrarea}

În planificarea urbană există conceptul de dezvoltare pe spațiu verde. Aceasta se referă 
la dezvoltarea pe un loc în care la momentul începerii lucrărilor este spațiu varde ce nu necesită
costuri adiționale pentru curățare. Atunci când vrem să construim un sistem bazat pe microservicii
pornind de la zero, apare acest fenomen. Însă de multe ori există un sistem existent ce ajunge
la anumite limitări ce fac dorința de trecere către microservicii. Acest lucru face referire la 
conceptul de dezvoltare pe spatiu maro, adică pe un spațiu post industrial în care există costuri de 
demolare sau restaurare și ce necesită o atenție mult mai ridicată datorită posibilității alterării
compoziției pământului ce ar putea să fie toxic.

Migrarea unei aplicații monolitice nu ar trebui să fie făcută cu scopul a obținerii unei aplicații
formate din microservicii ci pentru că există o funcționalitate ce este împiedicată de a
funcționa la capacitățile necesare pentru funcționare optimă, astfel migrarea unui sistem 
monolitic ar trebui să fie făcută incremental și posibil nu complet dacă nu este nevoie.

Începutul migrării ar trebui să fie făcut cu o analiză a sistemului și identificarea dacă 
avem posibilitatea de a întreține microserviciile, acest lucru înseamnă să avem prezentă infrastructură și
oameni care să știe să configureze uneltele necesare precum cele de CI/CD pentru a livra cu aceași consistență.
Ulterior putem analiza sistemul pentru a observa locurile unde se formează coeziune fapt
ce duce la posibilitatea creării unui microserviciu cu cod grupat. Pentru acest lucru se pot folosi 
unelte precum CodeScene ce pot identifica porțiunile de cod ce se schimbă împreună. Acest lucru 
ajută în crearea unui context mărginit și să identificăm modul în care acesta ar funcționa cu exteriorul.

O dată ce identificăm lucrurile ce vrem să le separăm putem separa funcționalități și elementele 
conținute, UI-ul, backend-ul și baza de date. Am fi tentați să ignorăm o parte din aceste componente
însă microserviciile funcționează cel mai bine atunci când toate acestea sunt separate prin simplul motiv
ca chiar daca UI-ul nu este 1:1 cu backend-ul, acestea tot vor fi modificate în același timp. În acest scop,
există două modalități în care putem face migrarea, pornind de la cod, însemnâand că separăm funcționalitatea
însă în continuare folosim datele din baza de date monolitică. Acest tip de abordare este destul de simplu,
întrucât obținem rezultate bune pe termen scurt, însă există limitări întrucât este posibil să nu
putem separa datele din baza de date, fapt ce înseamnă că munca depusă pentru pentru separare ar fi
inutilă. Separarea datelor mai întâi se referă la crearea unei baze date separate pentru datele 
pe care vrem să le folosim, acest tip de abordare aduce mai multe probleme la început întrucât 
ne face să ne gândim la probleme precum pierderea integrității datelor.

Atunci când facem modificări de tipul acesta trebuie să avem grijă cum orchestrăm schimbările.
O modalitate ce provine de la un arbust tropical „Strangler Fig” se referă la crearea unui strat
de intercepție ce verifică dacă funcționalitatea apelată a fost implementată într-un microserviciu,
și daca este atunci va ruta traficul către microserviciu. Acest tip de implementare permite să nu
modificăm aplicația monolitică și să putem să ne întoarcem la ea în cazul în care avem probleme.
De asemenea, dacă datele sunt comune putem folosi load balancing între cele două la nevoie.
În rularea paralelă, microserviciul și aplicația monolitică funcționează în paralel și ne permite să
le comparăm ca să ne asigurăm că avem funcționalitatea dorită. Metoda comutatorului ne permite
să avem cele două sisteme aplicația monolitică și microserviciul să funcționeze în același timp și să 
comutăm între cele două foarte ușor printr-un proxy.

În momentul în care începem separarea datelor, apar probleme ce nu pot fi prezise cauzate de lipsa
unei baze de date monolitice. De exemplu, performanța unei baze de date este extrem de bună când este vorba
de crearea de cereri ce conțin operații de alăturare de tabele, astfel noi trebuie să mutăm aceste operații
extrem de optimizate în interiorul microserviciilor. O dată cu lipsa unei baze de date comună pentru toate
tabelele apar problema lipselor de constrângeri, nu ne mai putem baza pe baza de date să 
forțeze aceste constrângeri și trebuie să verificăm și să tratăm erorile în cazul în care apar erori.
Tranzacțiile nu mai sunt garantate. O dată cu separarea datelor se pierde ACID-itatatea, o altă problemă 
a microserviciilor atunci când avem tranzacții distribuite între mai multe microservicii. 
În funcție de necesități ar fi necesar să rulăm cereri direct în baza de date sau să oferim acces
către persoane din exterior către aceasta, însă prin conceptul
de ascundere de informație acest lucru nu ar fi posibil. Ca să rezolvăm această problemă putem crea o copie
a bazei de date, read only, pe care microserviciul o va popula o data ce populeaza baza de date inițială.

Astfel, migrarea unui sistem monolitic este extrem de dificilă și vor apărea probleme, din acest
motiv este recomandat să folosim o abordare incrementală în care mutăm funcționalități una câte una,
acoperindu-ne obiectivul pe care ni-l dorim prin introducerea arhitecturii formate din microservicii.












\section{Tehnici de implementare}
Trecerea de la o aplicație monolitică la una bazată pe microservicii uneori nu este
ușoară. Unul din principalele motive este o posibilitate de a regândii modul cum accesăm
anumite elemente. Atunci când totul se află în cadrul unui singur proces, acessarea datelor
devine mai grea întrucât nu avem acces la ele și uneori nu este rentabil să căutăm după ele.
Atunci când apelăm o funcție local, timpul necesar pentru executarea acesteia este neglijabil,
atunci când această funcție trebuie să apeleze un alt microserviciu, lucrurile se complică și
trebuie decis dacă chiar avem nevoie. Un apel făcut printr-o rețea mereu va avea o întârziere, de asemenea
cresc locurile în care acesta poate eșua. În dezvoltarea aplicațiilor mobile întrucât 
ne așteptam ca aplicația să reacționeze la atingeri, nu putem să facem anumite operații pe
thread-ul principal, acestea pot fi apelarea lucrurilor de pe Internet sau accesarea unei baze
de date (chiar și dacă este de tipul SQLite3 fiind locală), aceasta pune în perspectivă 
impactul pe care acest tip de apeluri le pot avea.

Atunci când modificăm modul în care comunică un microserviciu trebuie să avem grijă la modul
cum ar afecta alte microservicii, sau atunci când acest lucru este necesar trebuie să orchestrăm
aplicarea acestor schimbări. Un alt concept ce într-o arhitectură monolitică se tratează ușor
este tratarea erorilor, întrucât toate se întâmplă în același mediu de lucru putem avea o viziune
asupra ce poate să cadă și în ce mod. Într-un sistem distribuit, pot apărea mai multe tipuri
de erori care se datorează rețelor, de exemplu crash-uri de microservicii ce necesită resetări,
erori de omisie de mesaje, când încercăm să apelăm un microserviciu însă acesta nu trimite un răspund,
în același timp acest răspund poate să fie mult prea rapid și să nu putem să îl recepționăm
sau prea târziu și să considerăm că nu am primit, dar și erori arbitrare precum unele în care 
ne așteptăm la un răspund și nu primim ce avem nevoie.

Toate aceste aspecte trebuie tratate atunci când lucrăm cu microservicii, însă fiind vorba
de o arhitectură paralelă ce oferă independență avem flexibilitatea de a alege ce folosim și ce nu.
De aceea este important să alegem tehnologii și stiluri de comunicare care oferă 
performanțele dorite încercând să minizăm timpul de răspuns și tipurile de avarii ce pot apărea.

\subsection{Tipuri de comunicare}
\section{Build}
Partea de construire a executabilului este o parte importantă din dezvoltare. Pentru limbajele
compilate este un moment în care putem să vedem dacă avem erori ce în general sunt destul de minore
însă sunt prinse încă de la începutul rulării ceea ce ne permite să le rezolvăm rapid, având
feedback aproape instant de la compilator.

În aplicațiile monolitice, crearea executabilului este destul de ușoară întrucât există 
un singur produs ce se află într-un singur loc ce trebuie să fie lansat, 
astfel avem un număr de executabil minim. În contextul microserviciilor, 
fiecare microserviciu are „executabilul” propriu, însă lucrurile se complică 
atunci când luăm in considerare că fiecare microserviciu are și o bază de date locală pentru el,
astfel aceasta trebuie să vină la pachet și să fie configurată și executată.

Un concept devenit standard este integrarea continua („Continious Integration”) - CI, aceasta
presupune ca atunci când lucrăm într-o echipă, codul cu care lucrăm să fie mereu merge-uit
într-un loc central pe care se rulează teste automate pentru a verifica că nu există probleme
la integrare, iar dacă există atunci acestea ar trebui să fie principala prioritate. Acest 
concept își are baza în Agile și partea în care vrem ca mereu produsul să fie într-o stare
în care poate fi livrat. CI-ul presupune să avem artefacte create din codul ce este merge-uit,
astfel beneficiem de un istoric al build-urilor ce ne permite să testăm pe versiuni anterioare
în cazul în care apar probleme.

Necesarul pentru un sistem de CI eficient este să integrăm codul frecvent, ideal zilnic,
pentru a putea valida că codul rulează bine împreună, acest lucru făcându-se prin teste automate
pentru a asigura că modificările aduse nu aduc schimbări de comportament, fiind validat 
prin teste dar și corectarea codului ce face ca testele să nu funcționează, acest lucru este necesar
întrucât cu cât mai multe modificări aducem unui produs într-o stare stricată, va fi mai dificil 
de depanat.

Un mod de dezvoltare a caracteristicilor noi este de a crea branch-uri separate pentru fiecare.
Dacă cerința pe care încercăm să o dezvoltăm necesită mult timp, asta înseamnă că nu o să integrăm
codul decât la finalizare, ceea ce nu este în corcodanță cu principiul integrării frecvente.
Astfel, e mai indicat să fie un singur branch în care cei care dezvoltă adaugă schimbări frecvent
ceea ce evită problemele de integrare, frecvent numit și „trunk-based development”, însă nu este
perfect întrucât modificările făcute de alții sunt mai puțin vizibile, întrucât pull-request-urile sunt 
făcute ca o atenționare asupra modificărilor care vor apărea în cod, din acest motiv
problema cea mai mare este să așteptăm primirea acceptului la pull-request-urile lansate.

Integrarea continuă se face prin intermediul pipeline-urilor care sunt configurate prin 
intermediul anumitor unelte care sunt făcute special ca să îndeplinească aceste cerințe,
în cadrul lor se rulează mai multe etape în care se face o parte din procesul de construire al
artefactului. Acestea sunt definite cu ajutorul aplicației de CI. Cel mai important pas
este cel în care se rulează testele. În general este bine ca testele care rulează repede
să fie executate la început întrucât cu cât feedback-ul pipeline-ului este mai rapid, cu atât
fixarea lui va fi mai rapidă. În cadrul pipeline-ului se crează un artefact asupra căruia
se vor executa restul pașilor.

Un alt aspect de considerat este modul de organizare al codului sursă, când avem o aplicație
monolitică aceasta nu are motiv să fie despărțită în mai multe repo-uri, însă în cadrul microserviciilor
avem opțiuni.

Cea mai simplă alegere este să nu facem nimic special și doar să punem toate microserviciile
într-un singur director, iar la final avem un singur build. Din punctul de vedere al unui operator,
acesta trebuie să lanseze cod și să configureze un singur loc, iar un dezvoltator indiferent
de cât de multe microservicii ar schimba, acesta face un singur commit. Însă o astfel de arhitectură
face lucrurile simple la început dar complicate la final. Având un singur build, timpul pentru feedback
crește întrucât se pot rula teste și pentru microservicii care nu sunt schimbate. De asemenea,
nu știm ce microservicii ar trebui să fie lansate și care nu întrucât la final ce ne face
să lansăm totul în același timp. Astfel, putem să începem cu o astfel de organizare însă
dacă vrem să avem un sistem de succes, eventual va trebui să avem o altă abordare.

Într-o arhitectură distribuită formată din mai multe programe, probabil primul gând ar fi să
avem o sursă pentru fiecare microserviciu, numit și multirepo este un mod de organizare
în care fiecare dintre acesta are spațiul și pipeline-ul propriu, astfel se produce 
o separare foarte ușoară între organizare și se cunosc echipele ce lucrează la fiecare
microserviciu. Însă simplitatea acestei separări aduce probleme în alte părți. Reutilizarea
codului devine mai complicată, întrucât fiecare microserviciu va depinde de librării
specializate, dacă microserviciul de la care pornește această librărie face o extensie 
atunci pentru a lansa caracteristicile noi trebuie să așteptăm ca microserviciile ce implementează
librărie să aducă o versiune actualizată. Arhitecturile distribuite contribuie la un tot unitar,
deci uneori o modificare trebuie să fie făcută în cadrul mai multor microservicii, în acest caz
trebuie să avem acces la ambele repo-uri, să facem commit-urile în ambele și după să actualizăm
și UI-ul ce consumă aceste microservicii. Dificultatea modificărilor crește mai ales dacă
e nevoie să fie implicați mai mulți oameni, însă dacă avem astfel de probleme e posibil 
ca modelul de microservicii să fie prea cuplat. Folosirea unei astfel de organizări 
merge suficient de bine indiferent de nivelul la care ne aflăm însă dacă facem modificări
des în mai multe locuri probabil ar trebui să facem câteva modificări la modul de desfășurare.

Un alt mod de organizare este de avea tot codul sursă într-un singur loc. Acesta sună asemănător
cu prima variantă prezentată însă aceasta este configurată și se cunoaște faptul că în cadrul
ei sunt mai multe proiecte. În primul rând, serverele de integrare continuă ar începe să diferențeze
între proiecte în funcție de poziția lor în director, astfel nu ar mai exista probleme cu
build-urile de lungă durată întrucât s-ar face per proiect modificat. Un avantaj al acestui
mod este faptul că reutilizarea codului este mult mai ușoară întrucât fișierele s-ar afla 
într-un loc foarte apropiat, singura dificultate ar fi să includem aceste fișiere în
serverele ce crează build-urile. Totuși separarea între echipe nu este la fel de evidentă,
pe GitHub există un tip special de fișier numit CODEOWNERS în care putem să specificăm 
pentru fiecare folder cine deține elementele, acest lucru ajută întrucât cei menționați
vor fi adăugați automat la recenzie atunci când se face pull request. Acest tip de pattern
are dezavantajul că ar consuma extrem de multă memorie, în același sens ar fi destul de greu
de descărcat mai ales când nu avem nevoie. În acest scop companiile mari își creează unelte
proprii pentru a rezolva aceste probleme, precum Google ce folosește unelte de versionare
a codului propriu, Piper, astfel acest tip arhitectural funcționează foarte bine pentru companiile
mari ce pot investi în îmbunătățirea acestuia dar și pentru echipe foarte mici în care 
separarea între proiecte este ușor de făcut.

\section{Testare} 
Testarea este o parte importantă în dezvoltarea produsului, avantajele creării de produse
sustenabile ce nu conțin defecte sunt destul de evidente întrucât clienții pot aprecia
faptul că aplicația nu se strică și pot construi o reputație în acest sens. Pe de altă parte,
apariția unei probleme majore în produs poate fi extrem de costisitoare în funcție de momentul
în care aceasta este găsită, astfel este important ca în procesul creării defectele să
fie prinse și rezolvate din timp.

Testarea poate să aibă o reputație scăzută fiind asociată ca ceva repetitiv ce trebuie făcut
manual însă aceasta a avansat și se face în diferite moduri inclusiv automat astfel poate
să fie la același nivel cu dezvoltarea însă necesitând un alt set de abilități și un mod
de gândire diferit.

Însă acest tip de testare se referă la oamenii dedicați pentru asigurarea calității produsului,
deci echipa de QA. Însă ei nu sunt singurii cu responsabilitatea de a monitoriza
calitatea sistemului, acesta pornește direct de la dezvoltator. În general, un tester poate
să nu aibă acces sau să nu se uite la cod, numit si black-box testing, în timp ce primind
acces se numește white-box testing iar atunci când avem acces parțial, gray-box testing.

Testele pot urmării lucruri diferite, astfel conform cărții „Agile Testing” de Lisa Crispin
și Janet Gregory putem împărții testele în funcție de ce încearcă să descopere. Testele de acceptare
testează dacă funcționalitatea este conform specificațiilor. Testarea exploratorie se face
în general manual și încearcă să verifice dacă există probleme în implementare. Testele unitare
sunt automate și verifică dacă nu s-a stricat funcționalitatea funcțiilor în urma unor modificări noi.
Testele de proprietate verifică performanța sistemului și capacitatea acestuia, se face
prin folosirea de unelte specializate și urmărește timpul de răspus, scalabilitatea sau reziliența.

Putem grupa aceste categorii și să obținem mai multe clasficări, testele de afaceri ajută
alte persoane să înțeleagă calitatea produsului (TU, TE) și testează o parte mai largă din produs,
testele ce vizează tehnologia sunt mai rapide și au componente care sunt automatizate (TU, TP).
Pentru a observa calitatea produselor se pot folosi teste care au ca scop asistența în dezvoltarea
pentru a nu apărea regresii (TU, TA) dar și unele care critică produsul în funcție de performanță
sau adaptibilitate (TE, TP). În general testele se fac în timpul dezvoltării, însă există teste 
care se fac pe produsul care este în piață.

Testele de serviciu au ca scop testarea unui microserviciu, acestea evită testarea UI-ului
ci doar a componentelor folosite de UI ca să aducă funcționalitate utilizatorului. Limitând
testele la un singur microserviciu, testăm capacitatea acestuia iar orice comunicare cu exteriorul
este mimată, astfel cauzele de eșec pot să fie strict legate de microserviciul testat
Aceste tipuri de teste, întrucât nu testează componenta grafică pot sa fie la fel de rapide
ca testele unitare însă întrucât pot intra în componență apeluri la baza de date sau la rețea.
Astfel în implementarea acestora există necesitatea înlocuirii comunicării cu exteriorul și se
face în general fie prin mimicarea datelor sau prin înlocuirea acestora. Mimicarea se referă
că la orice tip de apel al unei funcții, indiferent de datele de intrare, timpul la care a fost
făcut sau numărul de ori de câte este apelată acesta va returna un singur set de date. Înlocuirea
se face prin replicarea funcționalității astfel putem apela funcții care procesează datele de 
intrare și returnează date în funcție de acestea. Putem să dezvoltăm această idee și să 
încercăm să creăm servere de test care furnizează această înlocuire existând unelte pentru aceasta,
de exemplu mountebank ce poate să mimice date trimite prin HTTP/S, TCP sau SMTP însă nu poate să
mimice comunicarea prin evenimente.

Însă chiar dacă facem aceste teste, nu putem să garantăm că putem livra microserviciul.
Testele end-to-end se fac prin mimicarea utilizatorului și a diferitelor comportamente pe care 
acesta le-ar face, astfel verificând funcționalitatea sistemului prin interfața cu care el comunică. 
Acestea pot să implice multiple microservicii ce trebuie să nu aibă probleme, însă având o arie mai mare 
de acoperirea acestea sunt mai greu de testat iar atunci când pică nu se poate observa din 
prima posibila cauză. Una dintre problemele care arpare în aceste teste este nevoia de a rula
mai multe microservicii in același timp, astfel pentru fiecare din aceste microserviciu s-ar rula aceleași
teste. Acesta ar duce la o repetare de teste care nu s-au schimbat, astfel o metodă de a eficientiza
procesul este de a combina testele end-to-end pentru mai multe microservicii in același timp.

Aceste tipuri de teste acoperă o arie extinsă de cod iar fiecare dintre acestea poate să aibă
probleme, astfel pot exista în suita de teste unele teste care uneori pică deși uneori trec, cauza
unor astfel de situații poate să fie probleme temporare de rețea ce cauzează un timeout.
În astfel de cazuri ar trebui ca acest test să fie fixat, fie eliminat întrucât existența 
unui test neconcluziv poate duce la ignorarea acelui test, astfel este ca și cum acesta nu ar
rula de la bun început. 

Întrucât testele end-to-end acoperă mai multe microservicii apare problema asocierii 
responsabilității scrierii acestora. Întrucât în același timp mai multe echipe sunt implicate,
acestea toate ar avea o parte de responsabilitate însă scrierea unui astfel de test 
include cunoașterea sistemului întreg și observarea locurilor ce pot fi mai sensibile,
astfel companiile mari pot investii în crearea de echipe dedicate ce au ca rol scrierea, 
rularea și rezolvarea (când este cazul) a testelor end-to-end. Însă crearea unor teste
end-to-end complexe poate duce la timpuri mari de executare ale acestora, astfel 
feedback-ul compatibilității întârzie iar în acest timp se pot aduna mai multe caracteristici
adăugate de dezvoltatori, astfel se încalcă regula menținerii o prioritate a rezolvarilor
problemelor ce sunt asociate pipeline-urilor, în acest caz putem analiza suita de teste și să
păstrăm doar esențialele, mai ales când un build poate să pornească suite de teste end-to-end diferite
întrucât acesta este prezent în mai multe procese care se desfășoară în cadrul mai multor
microservicii.

Problemele cu testele end-to-end duc la crearea unor probleme asociate cu arhitectura paralelă,
întrucât acestea ar trebui să ruleze separat, atunci ar putea să fie și testabile independent,
astfel putem rula teste end-to-end mai rar dacă avem diferite caracteristici care să 
ajute la interoperabilitatea precum interfețele explicite sau prin folosirea unor unelte precum Pact 
putem să creăm o schemă a comunicațiilor exterioare ce include așteptările pe care le avem de la 
aceștia. Acestea pot fi rulate după de dezvoltatorii microserviciului și să cunoască 
dacă apar probleme de compatibilitate. Acestea sunt folositoare întrucât nu trebuie să mai rulăm
mai multe microservicii în același timp ci rulează ca testele unitare sau cele de serviciu în cadrul
unui singur microserviciu, astfel devin mult mai rapide.

Majoritatea testelor rulează în timpul perioadei de dezvoltare, însă putem rula teste 
și pe produsul final. Înainte de lansare se poate lansa o suită de teste numită Smoke, ce 
are ca scop verificarea funcțiilor de bază, de multe ori aceasta ar trebui să ruleze înainte de testarea
manuală întrucât căderea acestuia ar însemna lipsa funcționalităților de care am avea nevoie 
în mod normal, însă ca o ultimă măsură putem să îl rulăm alături de un test de Sanity ce conține
o gamă mai largă de de teste înainte unei posibile lansări. De asemenea, alte
teste în producție pot fi cele de A/B sau feature toggles sau Canary releases însă acestea
sunt mai mult tipuri de lansări de produs. Alte tipuri de teste pot fi cele care simulează
un utilizator normal în timp ce produsul este în piață pentru a verifica dacă funcționalitățile rămân 
la fel și nu există probleme de configurări ale mediilor de lucru.

În cazul apariției a unei probleme în producție, unele companii rulează simulări și rapoarte
a modului de lucru în cazul apariției unei probleme. În cazul acestora se pot definii concepte 
precum timpul mediu de reparare(MTTR) sau timpul mediu între erori(MTBF). Prima dintre ele
se poate rezolva ușor prin crearea unei restaurări a versiunii anterioare iar a doua prin 
îmbunătățirea capacității sistemului de a rezista la erori sau de a le trata. Aceste două metrici
sunt la fel de importante ca și testarea întrucât ambele au același scop, asigurare calității 
unui produs.

Majoritatea testelor rulate vor avea ca scop testarea caracteristicilor și validarea
că acestea funcționează în parametrii normali, însă datorită volatilității traficului
pot apărea scenarii complexe în care pot apărea un influx de trafic, astfel există 
dorința de a testa după funcționalitate, adica aspecte precum performanța sistemului.
În cadrul unei arhitecturi paralele aceasta poate cauza probleme întrucât într-o aplicație
monolitică comunicarea se face într-un singur loc, însă datorită necesității a comunicării 
colective pot apărea momente în care sunt necesare crearea mai multor cereri, acestea pot
cazuri în care pentru microservicii suprasolicitate să creze creșterea latenței sau
chiar acestea să nu mai poată fi accesibile. Astfel pot exista mai multe teste de performanță,
cele de enduranță ce testează sistemul pe o perioadă lungă de trafic dar în parametrii normali,
de stres în care testăm sistemul ăn afara parametrilor normali, ale spike-urilor ce sunt creșteri
mari in trafic pe perioade scurte de timp și de scalabilitate ce observă capacitatea 
sistemului de a servi atunci când numărul de utilizatori crește. Toate acestea ar trebui să 
se facă cunoscând funcționalitatea normală a sistemului pe trafic real, în absența acestei condiții
putem să ne concentrăm pe direcții pe care utilizatorul nu le ia. Într-o arhitectură paralelă,
la fel ca în securitate, capacitatea acestuia de a rula eficient este dată de capacitatea celei
mai încete piese astfel putem crea scenarii de test în care testăm lipsa unor componente ce
au ca scop creșterea rezilienței în cazul erorilor.


\section{Deployment}
Ciclul de viață al unei aplicații monolitice continuă să fie simplă și în contextul
lansării acesteia. În urma trecerii prin toate fazele pipe-ului de build și posibila
validare manuală acesta poate fi trimis către client sau urcat din diverse platforme
de hostare, posibil să existe o perioadă în care platforma nu ar putea să fie accesibilă
datorită necesității de a reporni, dacă nu folosim un proces mai amplu de lansare 
care ar acoperi aceste goluri. Dacă este implicată și o bază de date atunci aceasta poate 
să fie lansate înaintea aplicației și totul ar funcționa, dacă este invers atunci probabil
aplicația tratează cazurile în care aceasta nu este disponibilă.

În cazul microserviciilor, procesul este mai complex întrucât caracteristica specficiă 
acestora este capacitatea de a fi lansate și a funcționa independent de celelalte, chiar
dacă pot fi uneori limitate de absența celorlalte, astfel procesul menționat anterior 
se repetă pentru fiecare microserviciu. Totuși, fiecare microserviciu poate să fie de o tehnologie
diferită, poate ca baza de date a fiecăruia să aibă un anumit tip și anumite cerințe de lansare,
astfel flexibilitatea oferită de acestea se poate resimți atunci când este timpul lansării.

Până acum, toate conexiunile între microservicii erau făcute doar logic, nu cunoșteam 
detalii despre modul cum acestea vor fi aranjate sau unde vor fi lansate sau sub ce formă.
O caracteristică a microserviciilor este capacitatea de a scala un anumit serviciu ce simte
o cerere mai mare, în acest context putem avea mai multe instanțe al aceluiași microserviciu,
la fel cum am face și cu o aplicație monolitică în cazul în care am vrea să servim mai mulți
clienți, mai lansăm o copie a aplicației. În acest caz, comunicarea nu se va mai putea 
face direct către microservicii întrucât scalarea se face automat, ci este nevoie să comunicăm
cu un intermediar ce ar face rutarea către acestea care are și scop să echilibreze
traficul între cele două, astfel folosim un load balancer și vom comunica direct cu acesta 
și nu cu microserviciul propriu-zis. În cazul comunicării cu evenimente, trebuie doar să 
ne asigurăm că evenimentul este consumat de un singur microserviciu iar în urma se va 
lansa evenimentul răspuns.

Un factor important însă care a fost ignorat în lansarea microserviciilor, este baza de date.
Întrucât fiecare microserviciu trebuie să aibă propria bază de date pentru a-și ascunde starea
internă și să poată să aibă control asupra tuturor operațiilor ce ar avea loc pe 
domeniul său, microserviciul trebuie să fie lansat împreună cu aceasta independent. Însă 
în momentul în care avem mai multe microservicii de același tip, adică mai multe replici,
baza de date va fi defapt comună pentru toate microserviciile, adică toate replicile
vor folosi ca sursă de date același loc. În acest context apar probleme cu scalarea,
întrucât e posibil ca serviciul nostru să scaleze infinit, baza de date nu se va scala, 
astfel poate să devină un impediment în performanța sistemului, dacă nu cumva își cauzează
sieși atacuri de tipul Denial of Service prin existența prea multor cereri. Însă scalarea
bazelor de date este complicată, mai ales în cazul în care acestea sunt relaționale, întrucât
necesită fragmentarea datelor ceea ce este greu de calculat, din acest motiv putem 
construi organizări precum existența unei singure baze de date pentru scriere și multiple
baze de date pentru citire, în acest context putem separa traficul însă datele pe care le 
servim pot să nu fie actualizate. Modul de lansare al bazei de date poate să difere de cel
al microserviciilor, mai ales pentru companiile existente ce au deja date, migrarea devine
mult mai dificilă, din acest motiv aceștia pot preferă să folosească infrastructura existentă
în loc să migreze către cloud unde ar putea avea o bază de date per microserficiu, astfel microserviciile
pot avea baze de date logic diferite însă care rulează de pe aceleași server.

\section{Monitorizare}
Necesitatea monitorizării nu apare atunci când dezvoltăm microserviciile ci atunci când le lansăm.
Atunci când apare o problemă nereproductibilă, singurul lucru ce ne poate descrie modul în care am ajuns în
această situație sunt urmele pe care le creăm intenționat, adică fișierele de jurnalizare
sau tranzacțiile care s-au desfășurat. Într-o aplicație monilitică acesta are un număr mic 
de locuri în care poate să cedeze, astfel investigațiile pornesc mereu din același loc. Monitorizarea
acestuia este mai simpla întrucât numărul de resurse pe care acesta le apelează este mai mic.
În cazul unei arhitecturi distribuite atunci când apare o problemă, acleași investigații
pe care le-am face pe o platformă monolitică o să le facem asupra mai multor componente.
Din acest motiv trebuie să găsim metode de monitorizare eficiente ce sunt ușor de accesat
și folosit pentru cei implicați

Implementarea unui astfel de sistem în interiorul unui microserviciu depinde de modul în care
acesta comunică. În cazurile simple în care avem o singură instanță al unui microserviciu
atunci monitorizarea acestuia este simplă întrucât trebuie să colectăm date dintr-un singur loc.
Atunci când acest microserviciu este scalat și există mai multe instanțe al acestuia care
pot să apară sau să dispară în funcție de load, atunci trebuie implementat un sistem
care colectează datele de la acestea și eventual le asociază datele asupra cererilor. Atunci când
aceste microservicii comunică și cu exteriorul devine și mai complicat întrucât nu vom știi
exact modul în care sistemul va acționa, în acest caz agregarea informațiilor nu este suficientă
și avem nevoie de metode de prelucrare eficientă al acestora.

Observabilitatea este capacitatea de a întelege starea unui sistem în funcție de datele pe care
acesta le comunică, asemănător cu bordul unei mașini ce ar lumina anumite beculețe în funcție 
de problemele pe care mașina ar putea să le întâmpine în interior. Aceste date furnizate pot
ajută în prevenirea unui eventuale probleme, fie prin scalarea automată a sistemului în cazul în care
datele furnizate de acesta indică că este suprasolicitat fie prin identiciarea problemelor 
înainte să apară prin crearea de algoritmi bazați pe inteligența artificială, iar toate acestea
ar trebui să se vadă dintr-un loc central pentru a putea examina starea sistemului în același 
mod pentru toate componentele acestuia.

În acest scop, atunci când creăm un sistem bazat pe microservicii ar trebui să pornim de 
la construirea unei baze pentru observabilitate, astfel putem avea în vedere anumite 
caracteristici care trebuie să apară. Una dintre cele mai importante și simple de implementat
acțiuni este agregarea log-urilor. Dacă putem crea asta atunci putem să construim un istoric
al acțiuniilor ce s-au declanșat în cadrul sistemului. Deși există unelte dedicate, toate
ar lucra în același stil, microserviciile crează fișiere de jurnalizare local pe care alt componentă
le va colecta și le va trimite către componenta de agregare. Este important ca acesta să nu există procesări mari 
ale acestor linii întrucât cantiatea de informație care apare poate să creeze creșteri mari în
consumul de resurse. Fiecare intrare în aceste fișiere ar trebui să fie standardizată cu informațiile 
de care avem nevoie, de exemplu data, microserviciul, mașina, nivelul de importanță al log-ului și 
atle date de care am avea nevoie în depanare. Toate aceste date sunt folosite de agregator pentru 
a crea centralizări informațiilor sau pentru a cunoaște momentele în care ceva grav se întâmplă
și să fie nevoie să alertăm personalul. Întrucât microserviciile funcționează pentru rezolvarea
unei singuri cereri, este important să cunoaștem cererea pentru care se lucrează,
în general se asociază un ID de corelare pe care fiecare microserviciu îl va comunica 
pentru ca tranzacțiile din cadrul cererii să fie ușor de interogat atunci când acestea vor fi
agregate. Deși fiecare înregistrare are o dată, aceasta nu poate să fie sigură iar încrederea în 
acestea poate fi o capcană. În programare, cel puțin pentru mine, lucrul cu timpul este cel mai dificil,
în interfață deoarce există fusuri orare ce trebuie tratate (desigur e ușor să salvăm datele 
în UTC și doar să le afișăm în mod diferit, însă se adaugă și schimbările create de economiserea
de economiserea de lumină), iar în cadrul fișierelor de log ale microserviciilor 
timpul nu poate fi precis întrucât fiecare are timpul său intern, deși poate să fie sincronizat
cu un server de sincronizare Network Timing Protocol (NTP) nu poate să fie exact. Putem folosi 
unelte precum Fluentd ce trimite date în Elasticsearch cu interfață în Kibana însă acest spațiu
este populat cu soluții. Atunci când avem un sistem de acest tip, pot apărea probleme, pe lângă ce a
fost menționat, apar costuri cu privire la stocarea datelor dar și necesitatea unei puteri computaționale
în creștere pentru a menține index-urile coloanelore din bazele de date.

Un alt aspect ce comunică starea sistemului sunt metricile și telemetria. Datele colectate de 
acestea pot include procentul de resurse consumate, arhitectura sistemului, modul cum răspund la cereri 
sau chiar performanța acestora. Aceste metrici sunt în genera simple întrucât urmăresc 
stocarea unei singure informații sau acțiuni, însă acompaniate de acestea apar informații 
suplimentare iar acestea sunt costisitor de depozitat întrucât va trebui să fie interogate.
În acest scop putem folosi unelte precum Prometheus sau Graphite.

Colectarea datelor despre un sistem se face într-un singur scop, aprecierea dacă totul
funcționează în parametrii nominali. În cadrul unei aplicații monolitice, această apreciere 
era simplă întrucât lucrurile pot să meargă sau să nu meargă. Într-o arhitectură paralelă
ce se bazează pe independență această apreciere devine mai dificilă. Dacă un microserviciu rulează,
este un pas, însă nu este suficient. Acesta trebuie să fie accesibil pentru toți consumatorii 
acestuia. Însă dacă funcționarea acestui microserviciu depinde de alt microserviciu care 
momentan are probleme, putem să spunem că acest microserviciu este într-o stare bună?
În aceste cazuri se creează anumiți termeni, primul concept fiind Termeni la nivel de serviciu,
„Service level agreement” ce se crează între cei care construiesc un sistem și cei care folosesc
acest serviciu, astfel aceștia din urmă pot să își creeze un anumit grad de încredere sau
nivelul la care se așteaptă. De exemplu, providerii de infrastructură în cloud pot 
crea astfel de termeni pentru a garanta funcționalitatea serviciilor oferite de ei, pe
perioadele în care acestea nu pot fi accesibile din cauza lor, probabil nu vor exista credite 
consumate însă probleme ale sistemului nostru tot pot apărea ce pot afecta clienții proprii.
Obiectivele la nivel de serviciu („Service-Level Ojective) sunt create pentru a crea o imagine sau un posibil progres
în atingerea SLA-ului. Accestea pot include lucruri precum timp stabil sau latență. Pentru 
a aprecia atingerea SLO-ului putem crea Indiccatori la nivel de serviciu („Service-level Indicators”)
ce depind de sistemul oferit. În construirea unui sistem putem să alocăm spațiu pentru erori, 
acesta este scopul SLA-urilor, întrucât nu putem să fim perfecți ar trebui să păstrăm
loc și pentru apariția unor erori. Acestea ajută în crearea deciziilor care includ riscuri 
precum înceracare de tehnologii noi.

Apariția erorilor este inevitabilă, din acest motiv trebuie să întelegem modul în care
acestea pot să fie raportate. Există erori al căror apariție poate fi ignorată pe moment, în 
general la aceste tipuri de erori ne putem aștepta și au o severitate redusă. În interiorul unui sistem
trebuie să știm tipurile de erori care impactează funcționalitatea sistemului și trebuie remediate
în cel mai scurt timp, în general pentru cazurile acestea un tehnician ar fi alertat 
prin diferite forme pentru a remedia problema. Întrucât astfel de situații crează un grad 
mare de disconfort, trebuie să alegem erorile care generează acest tip alerte cu grijă,
întrucât dacă sunt prea multe e posibil ca acestea să fie ignorate. Studii create
pentru a studia motivele pentru care unele accidente apar pot să concluzioneze
că din cauza numărului mare de alerte raportate care nu au gradul corespunzător pot genera
pentru cei care le urmăresc o oboseală, din acest motiv aceștia pot să nu urmeze procedurile
corespunzătoare. Steven Shorrock menționează în articolul „Alarm Design: From Nuclear
Power to WebOps” (Humanistic Systems (blog), October 16, 2015) că motivul alertelor
este de a direcționa atenția utilizatorlui către aspecte ale operației ce au nevoie
de atenție într-un timp scurt. („The purpose of [alerts] is to direct the user's attention towards 
significant aspects of the operation or equipment that require timely attention.”). Din acest motiv
ne putem inspira din reguli vizează comportamente ce necesită utilaje complexe.
Asociația utilizatorilor de echipament și materiale (Engineering Equipment and Materials Users Association - EEMUA)
a creat o descriere pentru crearea alertelor într-un mod corect. Acestea trebuie să fie relevante,
să ilustreze o caracteristică importantă a sistemului, unice, pentru a evita repetiția, să ajungă 
într-un timp util în care acțiuniiile utilizatorilor mai pot crea un impact, sa conțină
suficiente informații pentru ca acesta să le poată prioritiza, informația ce provine din alerte
să fie clară și ușor de înteles, să conțină informațiile necesare pentru a începe o acțiune cu scopul
de fixare, dacă este posibil să ofere sfaturi pentru rezolvare, și să fie concentrate pe 
problema pe care o raportează.

Monitorizarea unui sistem este dificilă, însă aceasta poate să fie introdusă treptat însă 
putem începe cu agregare fișierelor de jurnalizare dar și notarea parametrilor în care rulează
mașinile de lucru și raportarea acestuia într-un loc central.

\section{Securizare}
Indiferent de platforma sau tipul dezvoltării pe care îl facem, front-end, back-end sau IoT,  
securitatea va fi una din problemele care vor fi ridicate atunci când dezvoltăm produsul.
Însă aprofundarea acestei părți nu este ușoară întrucât necesită o mentalitate specială în 
care luăm în considerare un adversar. În general, în timpul dezvoltării construim elemente 
și le construim într-un mod în care să funcționeze eficient, însă există persoane care 
vor să abuzeze de vulnebaribilățile ce pot apărea la orice nivel al aplicației.

Din acest motiv, în cadrul unei companii există persoane speciale care se ocupă de securitate.
Aceștia cercetează și studiază posibilele vulnerabilități care apar și crează anumite investigații
asupra produsele companiilor. În același sens, se pot folosi și unelte speciale în interiorul pipe-ului
care au ca scop scanarea fișierelor și raportarea posibilelor vulnerabilități, chiar dacă nu este amănunțit
sau o scanare în detaliu, aceasta ar oferi feedback rapid ce poate fi soluționat în același mod.

Prin stilul lor, microserviciile ne oferă posibilitatea de a îmbunătății securitatea sistemului
întrucât fragmentarea duce la ascunderea informației către serviciile din exterior, astfel un
atacator ar primi de fiecare dată un minim de informație, însă pe cealaltă parte, microserviciile sunt
foarte complexe și necesită mult mai multe setări, pe lângă infrastructură și datele pe care 
acestea le păstrează, cresc numărul de rețele ce trebuie securizate, numărul de dependințe
ce pot proveni de la terți (Eventbus, load-balancers, log tools) cresc și trebuie scanate
regulat pentru vulnerabilități, astfel numărul de sisteme ce trebuie să fie luate în considerare 
atunci când construim metodele de protejare ale sistemului.


\section{Evolutie}
\subsection{Stabilitate}
Într-un proces clasic de dezvoltare a software-ului în stilul Waterfall, după ce aplicația
este construită, testată și trece de audit-ul echipei de securitate aceasta va fi lansată și
se va ajunge în starea de mentenanță. Însă modern, dezvoltarea aplicației nu se termină
aici ci se va dori extinderea, consumul aplicației de către diferite interfețe, scăderea
numărului de erori și tratarea acestora dar și eficientizarea sistemului.

Fiind implicată conexiunea prin Internet dar și globalizarea, utilizatorii se așteaptă
ca serviciile pe care le oferim să fie disponibile 24/7. O arhitectură distribuită ne
permite să creștem această disponibilitate întrucât poate să se întâmplă ca un serviciu
pe care îl creăm să fie căzut, celelalte fie nu ar fi afectatate sau ar rula într-o variantă
mai limitată.

Pentru a crește reziliența sistemului putem definii anumite arii în care aceasta acționează.
Când aplicația poate să trateze erori pe care le anticipăm, precum tratarea erorilor sau
reîncercarea cererilor căzute. Din această cauză putem să construim soluții ce devin mai complicate,
de exemplu adoptarea Kubernetes pentru a beneficia de facilitățile acestuia precum
menținerea constantă a stării, introducem complexitate și un efort ce poate să creeze mai multe
probleme. Capacitatea de a absorbi atunci când se produce o anomalitate sporește stabilitatea
sistemului, concentrarea pe îmbunătățirea metodelor de prevenire al cazurilor de acest
tip ne poate lăsă vulenaribili, astfel putem crea metode de recuperare sau metode alternative
de funcționare atunci când apar probleme în sistem. Creșterea rezilienței ar trebui să
fie făcută constant și să ne adaptăm la evoluția sistemului.

Deși putem încerca să reducem șansele de apariție ale unei probleme, acestea nu pot fi evitate,
iar acestea pot apărea la orice nivel al aplicației inclusiv în afara ei la partea de infrastructură.
Însă tratarea tuturor evenimentelor care ar apărea ar fi obositor și nu ar da randament,
întrucât pot fi microservicii în sistemul nostru la care ne putem aștepta să cadă și să
nu creeze perturbări la fel de mari. Astfel putem să testăm monitorizăm sistemul și
să calculăm timpi de răspuns, dacă aceștia nu sunt în parametrii pe care îi dorim
putem folosi sisteme de scalare, sau a timpului în care sistemul este căzut și stabilirea
dacă aceasta este acceptabilă la fel și cât de multe date putem să pierdem.

Pentru a continua funcționalitatea sistemului în procent cât mai mare, uneori e nevoie
să facem compromisuri, astfel dacă cunoaștem că un microserviciu de care suntem
dependenți este căzut, putem să îl ascundem, astfel utilizatorul are o experiență mai proastă,
însă sistemul încă este disponibil. Metode de menținere a stabilității sistemului pentru
astfel de situații pot fi, introducerea de time-out-uri în cereri, astfel clienții nu așteaptă
mai mult decât trebuie. Durata după care se introduce un timeout poate fi calculată
după rularea testelor și adjustată conform telemetriei. Însă unele erori pot să fie
doar temporare, iar o reîncercare probabil ar fixa problema, în acest caz putem
introduce un sistem de reîncercare în funcție de eroare pe care o primim. Aceste două metode
de tratare a eșecului totuși sunt costisitoare din punctul de vedere al timpului. Întrucât
conexiunile încete și un număr ridicat de reîncercări ar face ca clientul să trebuiască
să aștepte foarte mult, chiar și pentru o eroare. În acest caz putem crea un fel de siguranță
ce are scopul ca atunci când avem funcționalitatea normală aceasta să continue, iar atunci
când avem probleme în sistem, clienții care ar încerca să acceseze resursele afectate ar primi
un răspuns cât mai rapid, astfel nu ar cauza frustrări pentru aceștia. Erorile se pot
propaga, astfel crearea unei izolări între microservicii ar cauza ca acestea își afectează
starea reciproc. În acest sens putem crea separări logice în care comunicare între microservicii
se face printr-un intermediar care în absența unuia dintre participanți va salva lucrurile
ce trebuie să le trimită și le va trimite atunci când va fi disponibil, sau izolare
fizică, astfel să ne asigurăm că bazele de date ale microserviciilor dar și aplicațiile în sine
sunt separate între ele, astfel dacă un server cade, nu ar afecta mai multe servicii în același timp.
Deși acest lucru este mai costisitor putem să folosim asta și să coordonăm metodele
de eșec, astfel dacă știm că funcționalitate unui serviciu este puternic afectată
de un altul, atunci acestea pot fi grupate și dacă inevitabilul se întâmplă cele două
vor cădea împreună. În același stil putem să creștem numărul de apariții ale unui microserviciu,
astfel dacă unul din ele este căzut, avem alte clone ce îi pot lua locul. Platforme precum
Kubernetes ne permite să facem acest lucru cu ușurință. În același stil, putem să trimitem
mai multe cereri de mai multe ori, acest lucru funcționează dacă cererile create și modul
lor de prelucrare este idempotent, adică de oricât de multe ori este aplicat, acesta
va da mereu același rezultat oricât de multe rulări am face.

În construirea unui sistem bazat pe microservicii stabil am discutat despre metode de
eficientizare și de îmbunătățire a rezilienței în care folosim mai multe baze de date sau
multiple instanțe ale unui microserviciu. Există o teoremă ce a fost demonstrată și matematic ce
spune că putem avea consistență, disponibilitate (avalability) și tolerare a separărilor
(partition tolerance) însă dintre acestea putem să alegem în sistemul nostru doar două dintre ele,
a treia neputând fi atinsă. Consistența reprezintă ca oricare din instanțele ale unui microserviciu,
toate ar oferi același răspuns pentru aceași cerere. Disponibilitatea reprezintă capacitatea
unui microserviciu de a primi un răspuns iar toleranța separării reprezintă capacitatea
sistemului de a manevra incapacitatea serviciilor de a comunica între ele.

Renunțând la consistență ar însemna ca atunci când avem mai multe microservicii de același tip ce comunică
cu o bază de dată, dacă nimerim la noduri diferite am primi răspuns diferit. Acest lucru se
poate întâmpla atunci când baza de date este formată din mai multe baze de date, unele
destinate doar citirii iar una doar pentru scriere. În acest context, sincronizarea
dintre cele două tipuri de baze de date este pe moment întreruptă. În general putem
sacrifica un pic din consistență în funcție de importanța datelor cu care lucrăm, însă
ne așteptăm ca sincornizarea să se facă la un moment dat, cu cât aceasta întârzie cu atât
va fi mai dificil de executat, iar dacă refuzăm scrieri atunci când sincronizarea
este lentă o să menținem consistența ridicată dar disponibilitatea scăzută.

Consistența este greu de atins întrucât necesită mulți pași pentru menținerea ei, iar interogarea
datelor se face inițiind tranzacții pentru a ne asigura că datele pe care vrem să
le accesăm sunt cu adevărat la fel, iar dacă aceste noduri nu sunt disponibile, atunci pentru a
nu sacrifica consistența vom refuza cererea, astfel scăzând disponibilitatea. Crearea unui mod
de partajare a datelor între mai multe noduri este dificilă, astfel ar fi bine să folosim
soluții deja implementate, Consul oferă astfel de servicii pentru disitribuirea elementelor
de configurare.

Negarea tolerării separări este imposibilă în cadrul unui sistem distribuit, întrucât aceasta
ar însemna să renunțăm la microservicii și să revenim la o arhitectură monoliticăm în care
nu există separare. Astfel singurele tipuri de sisteme pot fi CP sau AP, adică cele care
renunță la disponibilitate, respectiv la consistență, însă în lumea reală acesta nu
este neapărat aplicabil. Putem în cadrul unui singur sistem să avem mai multe tipuri de alegeri,
în funcție de datele și tehnologia cu care lucrăm.

Testarea stabilității unui sistem distribuit este destul de complicat și s-ar face la mai multe nivele.
Un termen popular, original provenit de la Netflix, este „chaos engineering” (ingineria haousului)
acesta are ca rol testarea rezilienței sistemului. În general, poate fi considerat
ca doar rularea unor unor unelte și interpretarea acestora, ceea ce nu e greșit, însă
aceasta are ca scop îmbunătățirea încrederii atunci când apar probleme ce nu pot fi
prevăzute, ceea ce se poate întâmpla în orice moment. Netflix se bazează într-o proporție
foarte mare de infrastructura ce provine de la AWS, însă la un moment dat aceștia au fost
forțați să efectueze lucrări de mentenanță asupra unei întregi zone de serviciu, astfel
se putea crea o gaură mare de down-time pentru clienții Netflix (dar și pentru ceilalți ce se
bazează pe AWS). Ei au creat un serviciu numit Chaos Monkey ce are ca scop testarea sistemului
la apariția unor probleme neprevăzute în cadrul infrastructurii. În cadrul serviciului
putem întâlnii mai multe stagii precum, Chaos Gorilla ce are ca scop căderea unei zone de disponibilitate,
adică a unei regiuni mai mici (Europe Central), Chaos Kong, simulează căderea unei regiuni întregi,
de exemplu Europa, Latency Monkey are ca scop creșterea latenței în cererile executate
de clienți pentru a simula încetiniri în trafic, Conformity Monkey are ca scop
închiderea instanțelor ce nu sunt în conformitate cu practicile cerute de sistem,
Doctor Monkey monitorizează starea fiecărei instanțe și oprește sistemele ce nu sunt tratate în timp util,
Janitor Monkey observă infrastructura ce nu este utilizată și o oprește după un anumit timp,
Security Monkey are ca scop închiderea serverelor ce nu au toate practicile de securitate
necesare.

\subsection{Scalare}

În timpul prezentării microserviciilor, un avantaj al acestora a fost mereu flexibilitatea,
în același stil am menționat posibilitatea îmbunătățirii sistemului prin replicarea unor
anumite funcții ale acestuia, fie pentru administrarea a mai mult trafic, fie pentru
îmbunătățirea disponibilității și eliminarea timpilor în care serviciul este inaccesibil.
Scalarea poate fi făcută în mai multe moduri, în funcție de necesitățile și de problemele
pe care le-am rezolva.

Scalarea verticală se referă la îmbunătățirea sistemului de calcul
prin aducerea mai multor resurse de calcul, fie prin îmbunătățirea sistemelor de I/O,
adică folosirea unor SSD-uri rapide față de HDD-uri clasice, fie de calcul prin
îmbunătățirea procesorului și a memoriei RAM. În cazurile clasice, când noi suntem
cei care hostăm aplicația, aceasta poate fi mai de lungă durată întrucât sunt mai multe
echipe implicate, însă având în vedere sistemele cloud, scalarea verticală se face foarte rapid
și nu necesită configurări adiționale. Însă acest tip de scalare nu este o soluție ce va
funcționa la nesfârșit, se observă că legea lui Moore nu se mai aplică, astfel îmbunătățirile
tehnologice nu sunt la fel de drastice în fiecare an. Nu ar trebui să ne bazăm că putem
schimba mașina de calcul, mai ales că folosirea tehnologiilor foarte puternice ar deveni
costisitor. În acest scop, putem folosi algoritmi mai eficienți sau alte tipuri de scalare.

Scalarea orizontală se referă la multiplicarea instanțelor ale unui serviciu, în acest
stil putem crește numărul de procesări ce se pot face în același timp. Implementarea
cea mai simplă se poate face folosind un balansator de încărcătură (load balancer) ce va
dirija traficul asigurând că traficul este egal pe fiecare dintre microservicii. Acest lucru
este posibil atunci când putem să distribuim o acțiune către mai multe locuri,
însă putem avea situația în care distribuim sarcinile din interiorul unei cereri,
astfel o altă implementare ar fi separarea unor cerințe din interiorul unei cereri
către mai multe microservicii ce au unic, de procesare, iar scalarea s-ar face asupra acestora.
Am mai menționat și scalarea bazei de date prin folosirea de replici de citire ce sunt sincornizate
cu baza de date destinată scrierilor. Limitările acestui tip de scalare se rezumă la dificultățile de
implementare, astfel putem întâmpina mai multe probleme de rețelistică sau de sincronizare a datelor
sau chiar recrearea unui nou mod de desfășurare destinat acestui stil distribuit de lucru.

Partiționarea datelor se referă la distribuirea cererilor către un anumit set de microservicii
în funcție de natura cererii. Astfel, putem avea microservicii de același fel, cu funcționalitate
identică însă care se ocupă de date diferite. Aceasta poate să fie o îmbunătățire întrucât
am avea sisteme de procesare dedicate unor anumite cereri, astfel dacă am întâlnii mai mult
trafic într-un anumit loc, am putea planifica o partiționare din nou, iar dacă nu avem
atunci să regrupăm anumite date. Implementarea se poate face în funcție de tehnologie, astfel,
dacă baza de date pe care noi o folosim oferă asistență în acest loc (de exemplu Cassandra)
atunci putem să o facem la nivelul acesteia, sau să implementăm servicii de proxy ce redirecționează
traficul către microserviciile ce ne oferă funcționalitatea de care avem nevoie. Acest tip
de scalare este eficient când avem suntem constrânși din punct de vedere al scrierilor,
întrucât acestea nu pot fi eficient distribuite în alt mod. Limitările pot fi
faptul că nu se aduce o îmbunătățire la fel de mare și în funcție de modul în care ne facem
grupările ne poate duce să avem un grup de persoane afectat, astfel putem grupa
acest tip de scalare cu cele precedente pentru efectivitate mai bună. În același stil,
stabilirea modului de grupare și dimensiunile partițiilor pot fi dificil de stabilit.

Majoritatea formelor de scalare încercau să nu modifice structura aplicației prea mult,
însă după ce epuizăm aceste forme, putem încerca extinderea aplicației prin decuplarea
funcționalității în alte microsisteme. Beneficiul acestui tip ar fi posibila eficientizare
a procesului întrucât s-ar putea parta încărcătura însă ar crește complexitatea sistemului
și poate cauza degradări ale perfomanței întrucât s-ar introduce o cerere pe rețea nouă,
astfel trebuie să ne asigurăm că implementarea acestui tip de funcționalitate ar avea
sens.

Sunt multe cazuri în care soluțiile pe care le creăm sunt mult mai complicate decât ceea ce
ne cere sistemul. Toate implementările de scalare sunt incluse în această frază, astfel
nu ar trebui să adăugam complixate fără să fie dovedit că am avea nevoie. În general,
există multe caracteristici introduse de dezvoltatori însă care nu sunt folosite mai niciodată
de către utilizator, însă atunci când aceste funcționalități au șansa să strice aplicația
ar trebui să fie evitate și aplicate doar la nevoie. Iar în acest context, întrucât microserviciile
ne oferă o flexibilitate ridicată putem să aplicăm la un anumit nivel toate tipurile de scalare
în loc să alegem doar una din ele.

Există și conceptul de auto scalare, în general ne putem gândii ca sistemul să fie mai puternic
accesat într-un anumit interval, astfel putem să creștem instasnțele de microservicii
în această perioadă. Ofertanții de infrastructură cloud ne oferă diferite metode
de autoscalare deja în platformă și sub diferite forme, de exemplu atunci când ajungem
la un anumită încărcătură pe un anumit nod pe o anumită perioadă putem să mai creăm o
instanță automat. Astfel autoscalarea poate fi reactivă sau predictivă, ambele aducând
îmbunătățiri însă trebuie să analizăm datele, ce poate fi făcut cu algoritmi complecși
precum cei de Inteligența Artificială / ML.

\subsection{Caching}

Caching este o metodă de a salva informația cerută la un moment dat de către un utilizator
și prevenirea recalculării modului de a ajunge la date și doar servirea aceluiași răspuns.
Atunci când putem să oferim același răspund întrucât l-am salvat putem spune că avem
un „cache hit” și când nu, un „cache miss”

Există multe motive pentru care am adopta această tactică, întrucât sistemul de cachare
are ca rost prevenirea recalculărilor inutile, putem îmbunătății performanța sistemului
întrucât unele dintre aceste calcule pot include cereri pe Internet ce adaugă latență.
În același stil putem eficientiza traficul dacă avem cachate anumite răspunsuri
pe care doar le livrăm de mai multe ori, astfel întrucât nu aduc complexitate pot fi
scalate cu ușurință. Deși ar scădea consistența, existența unui sistem de cache ne-ar putea
permite să rulăm o anumită perioadă de timp fără ca sursa obținerii datelor să fie disponibilă,
în acest stil putem ascunde problemele de puțină durată întrucât avem deja o sursă de date
introdusă.

Această tehnică poate fi aplicată la mai multe nivele, la nivelul clientului, putem salva
în interfață datele, de exemplu pe telefoane mobile putem face o dată o cerere către baza
de date, salvăm rezultatul și acesta va fi disponibil chiar și atunci când nu avem Internet,
acest lucru este vizibil în rețelele de socializare care atunci când deschidem aplicația
mereu au anumite postări preîncărcate, această interacțiune ne permite să salvăm baterie
și să nu creăm prea multe cereri, în același stil conținutul pe care îl afișăm utilizatorului
este mai greu de controlat și poate diferi de la unul la altul. Cacharea la nivel de server
ne permite să evităm anumite operațiunii de lungă durată pe care ne așteptăm să nu se întâmple
foarte des, astfel putem să trimitem un răspuns salvat. Această operațiune ne permite să
avem un control mai mare asupra datelor pe care le salvăm, însă clientul tot ar trebui să consume
resurse făcând cererea. În acest stil putem să salvăm informații ale unei cereri specifice,
iar dacă aceasta este repetată atunci se va trimite răspunsul vechi.

Atunci când vorbim despre salvarea unei informații prin caching, trebuie să luam în
vedere modul în care ștergem această salvare. Dacă nu am luat asta în considerare atunci clientul
ar vedea informații destul de vechi ce nu ar fi actualizate. Cel mai simplu mod de implementare
este specificarea unui timp la care informația din cache expiră. Intervalul poate fi dificil
de ales iar în același timp pot exista cazuri în care datele se actualizează fix după ce
clientul le înregistrează ca find cele mai recente iar în cazul în care TTL (Time To Live) este
setat mare ar putea cauza probleme. În același fel putem eticheta fiecare cerere cu un ID,
iar dacă ID-ul nu se schimbă atunci putem trimite un răspuns mult mai mic. Dacă avem o resursă
ce este disponibilă la un anumit URI, dacă am face cerere pentru aceași resursă putem trimite
un identificator al datelor pe care le avem iar server-ul ne poate trimite fie date noi,
fie faptul că datele nu s-au schimbat. O altă metodă de invalidare a cache-ului poate
fi folosirea de notificări, astfel server-ul va trimite fiecărui client care a cache-uit
un anumit mesaj, fie cu date noi fie doar faptul că au fost invalidate. Aceasta implementare
ar aduce avantajul că nu ar exista date vechi pentru clienți însă implementarea unui astfel
de mecanism este complexă și poate consumă destul de multe resurse. De asemenea, cache-ul
poate fi modificat fix după ce resursa se schimbă, o astfel de abordare ar avea sens
dacă am face cache-ing pe un server, sau să îl folosim ca un buffer în care scriem mai întâi
în cache și după acesta se va sincroniza cu serverul.

Caching-ul este o metodă de optimizare a traficului, astfel ar trebui să fie folosit
doar acolo unde am considera că ar crea o îmbunătățire a sistemului, la fel ca la scalare
acesta poate fi introdus treptat, deși poate face parte încă din arhitectura originală
însă putem să învestim resurse în ceva în care nu ar fi necesar și am aduce complexitate
într-un loc în care poate nu e nevoie.

\section{Organizare}
\subsection{Consum}

Atunci când construim un sistem, ne așteptăm să îl dăm în folosință pentru cineva,
fie clienți sau ca asistent al unui sistem existent. Însă în general acești clienți pot să
nu consume chiar sistemul direct întrucât acesta poate fi complex, ci putem să îi oferim printr-o
interfață construită tot de noi (în cazul aplicațiilor, putem considera front-end-ul o astfel
de interfață), astfel putem avea ca responsabilitate și crearea unui astfel de lucru.

În modul tradițional de lucru, am avea dezvoltatori grupați în funcție de rolul și platforma pe care
lucrează, astfel toți cei responsabili de o anumită componentă mai mare (de exemplu domeniu
sau prezentare) ar fi în același loc. Acest tip de organizare ne permite să avem o 
experiență identică la același nivel, însă comunicările între nivele pot fi afectate, astfel 
s-ar putea introduce probleme de sincornizare.

O metodă mai eficientă de organizare este gruparea dezvoltatorilor în funcție de locul în care
aceștia lucrează, astfel toți cei responsabili la toate nivelele de o anumită componentă a
sistemului nostru vor lucra împreună. Acest mod ne permite un control mult mai bun, având
metode de schimbări la orice nivel, astfel atunci când avem o inițiativă, aceasta e mai ușor 
de implementat și nu depindem de echipe exterioare pentru care trebuie șă așteptăm 
o implementare.

Construirea unei interfețe grafice poate devenii dificilă. În cadrul unei companii
ce oferă mai multe produse putem fi tentați să creăm o echipă dedicată pentru aplicațiile
pe care utilizatorii le-ar folosi, astfel aceasta ar fi consistentă tuturor. Această abordare poate
crea probleme atunci când vrem să lansăm caracteristici noi, de aceea o idee mai bună ar fi
crearea unei echipe ce va lucra în cadrul mai multor echipe, având posibilitatea să schimbe
proiectele.

Organizarea construirii interfețelor grafice poate fi făcută în mai multe moduri.
Cea mai simplă variantă presupune existența unei aplicații monolitice, în cadrul
acesteia consumăm API-urile propriu zise fără a face vreo optimizare modului cum
sunt făcute. Dezavantajul unei astfel de abordări este decuplarea între backend și frontend
dezvoltarea caracteristiciilor acestora nu poate fi la fel de bine coordonată. În același stil,
crearea multor cereri poate să cauzeze probleme legate de memorie, astfel acestea nu sunt
efectuate eficient.

O versiune mai avansată bazată pe framework-urile noi de dezvoltare este microfront-end-uri,
în cazul acesteia echipele furnizează un UI ce poate fi lansat independent de UI-urile
existente. O astfel de abordare le permite echipelor să își creeze UI-ul independent de celelalte
și acestea rămânând să se lanseze aproape independent, dar uneori tot combinat.

Această abordare poate fi implementată în mai multe moduri. Una dintre acestea este decompoziția
site-ului pe pagini, fiecare echipă fiind responsabilă de un anumit endpoint de pagini pe
site-ul mare. Aceasta le permite să fie independente întrucât livrarea paginilor este
independentă de paginile existente pe care server-ul trebuie să le arate.

O abordare mai dificlă de implementat, este separarea unei pagini în funcție de secțiile
pe care vrem ca acestea să le cuprindă, astfel fiecare echipă ar fi responsabilă de o componentă
ce apare pe o pagină la un anumit moment. Implementarea unei astfel de soluții ar fi ca
fiecare echipă să își construiască elementele ce vor să fie disponibile pe website și crearea
unei metode de agregare ale tuturor elementelor create. Din acest motiv, deși elementele sunt
separate acestea tot trebuie să fie integrate, iar crearea unui comportament împărțit
în cadrul mai multor elemente poate fi dificilă.

Unul din dezavantajele a unui front-end monolitic era numărul mare de cereri ce ar fi trebuit să
fie făcute către mai multe microservicii pentru o singură pagină. Pentru a rezolva problema aceasta
putem să folosim un intermediar ce are ca rol această agregare a datelor, astfel de pe dispozitiv
facem un singur call însă în spate sunt făcute mai multe. Dificultățile implementarii ar consta în
coordonarea echipelor, întrucât acest intermediar nu ar avea o echipă directă ci toți ar contribui la el.
În același stil ar exista dificultăți de tratare al numărului de informații pe care tipuri
de dispozitive le-ar folosi dar și de autentificare, întrucât totul devine agregat de un intermediar ce
ar trebui să aibă acces la toate resursele pe care le oferim.

Oferirea de asistență pentru diferitele platforme ne-a dus la crearea conceptului de Backend for Frontend,
astfel pentru fiecare tip de dispozitiv ce ar avea nevoie să ne acceseze sistemul,
am crea un backend specializat ce s-ar ocupa doar de o anumită experiență. Această abordare poate
avea sens dacă se fac multe apeluri peste rețea și dacă echipele de front end sunt decuplate de cele de backend.
Implementarea poate fi făcută de exemplu folosind GraphQL.

\subsection{Structură}

Nu este ceva nou faptul că aplicațiile mari sunt sunt create de mai multe persoane. Acestea pot fi
organizate în mai multe feluri și în același timp, fac parte din mai multe echipe și dintr-o
organizație ierarhică în cadrul companiei. Toate acestea au efect asupra modului în care
software-ul este dezvoltat.

În „How Do Committees Invent” publicat în Datamation în Aprilie 1968, Melvin Conway afirmă
că „Orice organizație ce crează un sistem, inevitabil va crea un design în care structura
va copia structura de comunicare a companiei” („Any organization that designs a system (defined more broadly here than
just information systems) will inevitably produce a design whose structure is a copy of the organization's
communication structure.”). Această „lege” practic afirăm că dacă am avea mai multe echipe ce ar
trebui să se ocupe de un sistem, atunci acestea își vor repartiza părțile astfe încât să lucreze fiecare,
deci se va ajunge asemănător la structura companiei.

Modul de organizare în companie dar și al echipelor influențează evoluția sistemului. Astfel,
unul din obiectivele pe care le putem crea este crearea independenței echipelor, acesta duce
la crearea unor organizații decuplate ce poate duce la eficientizarea modului de lucru.

Există probleme al căror eficientizare se poate face prin creșterea numărului de persoane,
în acest caz există multe exemple atunci când era prezentată regula de trei simplă, precum 
muncitorii pe șantier, aceștia au anumite cerințe ce pot să se desfășoare în paralel, astfel
adăugarea de mai multe persoane ar apropia data finalizării. Însă unele probleme nu se 
rezolvă în același stil, în absența capacității de paralelizare a task-urilor s-ar putea să 
înrăutățim situația. 

În acest context, alegerea unui număr potrivit de persoane într-o echipă poate fi dificlă.
Aceasta trebuie să aibă suficiente persoane în cât să fie independente. În cadrul Amazon,
s-a format conceptul de „Two-Pizza Team”, astfel în mod colocvial o echipă ar trebui să poată
fi suficient de mare încât să nu exite probleme în împărtășirea a două pizza în cadrul lor, 
și să fie suficient. Comunicarea ar trebui să fie eficientă în cadrul echipei iar între echipe
să fie limitată pentru a se preveni co-dependența în cazuri mari pentru că scopul 
este crearea unor echipe autonome.

Crearea unor echipe autonome duce la crearea unor deținători asupra microserviciilor create
de aceștia. Aceasta se numește o deținere puternică asupra acestuia, astfel doar echipa 
care deține microserviciul poate să dețină ce elemente intră în cod. Acest lucru se face
prin intermediul pull request-urilor, avantajul fiind crearea unui microserviciu stabil 
întrucât singurele modificări aduse sunt evaluate de persoane ce cunosc foarte bine proiectul.

În mod opus, putem avea microservicii ce au scop comun, astfel nu sunt deținute de nimeni și 
oricine poate aduce modificări acestuia. Acest mod de organizare ne poate aduce posibilitatea
să aducem modificări mult mai rapide, însă necesită ca echipele ce lucrează asupra acestuia
să se coordoneze mai bine, mai ales când intră în discuție tehnologiile utilizate.

Scopul final al organizării este de a eficientiza procesul de dezvoltare. Pentru a face asta, 
pot exista situații ce pot necesita acțiuni hibride.

De exemplu, combinând cele două abordări, putem avea un microserviciu deținut de mai multe echipe.
Acest lucru se poate întâmăla atunci când coeziunea este ridicată și decuplarea poate cauza
mai multe probleme decât ne dorim. Însă având mai multe echipe se pot cauza probleme în modurile de lansare
și introducerea schimbărilor ce pot strica anumite părți din alte microservicii, astfel 
aducem câteva din probleme unei arhitecturi monolitice însă la un nivel mai mic.

Pentru a crea o colaborare mai eficientă asupra codului comun, putem crea un mediu de lucru in 
stilul open-source intern, astfel toate echipele pot contribuii însă există anumite persoane mai implicate
în proiect care pot să aprobe shimbările. Acest mod poate avea sens și poate avea destul de eficient 
pe termen lung însă ar trebui ca proiectul să aibă un anumite de nivel de maturitate, astfel 
încât să existe o anumită arhitectură ce trebuie să fie urmată.

Un alt concept ce poate apărea este posibilitatea creării unui microserviciu dintr-unul existent.
Astfel putem avea o bază ce este liberă pentru toți și pentru funcționalitățile specifice 
se poate crea un synced fork al acestuia, asemănător moștenirii din OOP. Însă adoptarea unor 
astfel de schimbări trebuie să fie suportată de microserviciile cu care comunică.


\chapter{Orchestrare și infrastructură}
\label{orchestrare}
\section{Concepte generale}
Modurile în care un microserviciu poate fi lansat este destul de important întrucât
decide o bună parte din arhitectura inițială a acestuia, performanța sau capacitățile
pot fi eficientizate prin folosirea unei metode de deploy bună. Evoluția modului de a face rost
de infrastructură a pornit de la achiziționarea și creare echipelor special dedicate ce poate
devenii extrem de costisitor și necesitând mai mult timp, iar acum se poate face
adăugând cardul pe platforme specializate, adăugarea unor metode de configurare și
de securitate și putem avea un server într-un timp foarte scurt a cărui mentenanță la nivel
fizic devine responsabilitatea distribuitorului. Bill Baker, un inginer de la Microsoft
afirmă că la început tratam serverele ca și animale de companie, „le numim, și atunci
când devin bolnave, le îngrijești și le aduci înapoi. Acum ne tratăm serverele le tratăm
ca o ciradă. Le numeri iar atunci când se îmbolnăvesc, le împuști.” (You name them and when they get sick,
you nurse them back to health. [Now] servers are [treated] like cattle. You number them and when they
get sick, you shoot them.). Tehnologiile actuale ne permit să nu avem nevoie să încercăm să
reconfigurăm ci dacă putem salva datele de care avem nevoie, atunci este mult mai ușor să
reconstruim sistemul de la 0 și să îl aducem în starea în care avem nevoie de el.

Pentru a accelera crearea unui produs, nu este suficient să accelerăm dezvoltarea, ci găsirea
unei metode de eficientizare a întregului proces. În trecut, lansarea unui produs era una dintre cele
mai mari încetiniri întrucât includea crearea unui tichet și așteptara ca un specialist
să ne contacteze care ar primi programul și l-ar trimite în execuție. Procesul ar fi anevoios
întrucât nu am putea să lansăm atunci când am vrea ci doar atunci când ni se acordă timp, iar
în funcție de personal aceasta poate dura chiar luni, astfel lansările ar trebui să planificate
din timp și să fie solide, să nu existe necesitatatea de EFX-uri întrucât ar complica organizarea.

Cunoașterea metodelor de livrare și alegerea celei potrivite este esențială întrucât în
funcție de obiectivele pe care le avem, buget și alte considerente putem ajunge ca deployment-ul
să fie automat sau controlat la anumite nivele, astfel livrarea ar putea ajunge nici să nu mai fie
un punct dificil.

\section{Tipuri de deployment}
\subsection{Mașină fizică}
Cel mai simplu mod din punct de vedere conceptual, achiziționarea unui sistem de calcul,
configurarea acestuia pentru a avea un sistem de operare, crearea unor conexiuni
cu rolul de a ne face sistemul accesibil către locurile în care ar avea nevoie, descărcarea
programului pe care vrem să îl rulăm și executarea acestuia.

Această metodă are foarte multe constrângeri pe care în lumea microserviicilor sunt greu de
tolerat. Întrucât pe această mașină rulăm o singură aplicație, resursele consumate
vor fi doar provenite de aici, dacă sistemul nu este utilizat atunci am putea considera
că pierdem resurse.

În același stil nu avem metode ușoare de scalare, întrucât putem vedea că hardware-ul ales
nu este potrivit pentru aplicație și am dori să le schimbăm, acest lucru este ceva complex
ce necesită reconfigurări, migrări și cmoplexități de up-time ce sunt sezibabile.

Întrucât noi suntem cei care furnizăm platforma de calcul, atunci tot noi trebuie să îi
facem mentenață, acestea incluzând actualizări ale sistemului de operare, back-up de date
și audit-uri de securitatea, iar acestea pot consuma mult timp.

Deoarece sunt așa de multe limitări cu acest mod, în general preferăm să alegem unul din celelalte,
însă acest mod de deployment va fi mereu valabil pentru cazurile în care este nevoie de el.

\subsection{Mașină virtuală}

Întrucât lansarea direct pe hardware poate fi considerată costisitoare, am putea încerca
partajarea resurselor disponibile putând să rulăm mai multe microservicii pe aceași
mașină. Dacă am face acest lucru direct, atunci am încălca multe din bunele practici menționate
anterior, printre care izolarea rulării și lansării, adăugând complexitate configurăriilor.

Virtualizarea ne permite să eficientizăm resursele prin partajarea acestora între mai mulți
utilizatori. Aceasta este făcută prin crearea unor unelte ce încearcă să mimice elemtenele
unui calculator pentru fiecare mașină virtuală, astfel aceasta nu face diferența între
ce e fizic și ce e virtual și nu necesită o rescriere a unui sistem de operare pentru a-l
face compatibil cu virutalizarea, astfel acestea sunt independente între ele dar și de sistemul
gazdă putând să ruleze aproximativ orice.

Fiecare mașină virtuală poate să își configureze spațiul pentru microserviciul deținut,
acestea pot rula concomitent, astfel cantitate de resurse care este irosită este minimizată.
Aceste mașini virtuale pot fi limitate și le putem extinde în funcție de necesități. însă
crearea mai multor instanțe în cadrul aceleieași mașini nu va extinde posibilitățile fizice ale
acestuia și putem să ajungem într-un caz în care pe aceași mașină rulează prea multe mașini virtuale.
În lumea microserviciilor, astfel de grupări ar trebui evitate întrucât dacă mașina fizică
ar cădea, atunci tot ce e în interior ar fi căzut, însă dacă ar fi mai grupate am putea
crește viteza de răspuns dintre ele.

O altă considerentă este modul în care facem rost de mașini virtuale și modul în care acestea
sunt configurate. Furnizorii de infrastructură ne oferă API-uri detaliate, astfel sunt ușor
de adoptat cerințelor noastre, de asemenea ne oferă facilități de auto-scalare, back-up
și recuperare pe care în mod tradițional ar trebui să le configurăm. În mod clasic, companiile
pot folosi platforme precum VM Ware care deși îndeplinesc nevoile din punctul de vedere
al API-ului pe care il oferă poate fi considerat limitat.

\subsection{Container}

Problema mașinilor virtualizate era că atunci când le cream, fiecare dintre
acestea devenea o mașină independentă, asta ducea la duplicare de resurse, spațiu și procesare,
necesare pentru rularea sistemului de operare și ulterior a aplicației. Containerele
încearcă să rezolve această problemă prin partajarea mai multor lucruri, nu doar resursele
fizice dar și a nucelului, astfel în cadrul unui container putem rula orice sistem de operare
însă acesta rulează ca parte a ceva comun.

Procesele care rulează în container sunt suficient de bine izolate, acestea nu pot fi accesate
din exterior și nu pot comunica cu acesta decât dacă le permitem, aceasta se aplică și între ele,
sistemul de operare pe care rulează văzându-le ca niște simple procese.

Elementul care a continuat evoluția containerelor este Docker, acesta aduce beneficiile containerelor
însă aduce contribuțiile proprii asupra modul acestora de organizare. Acesta ne oferă asistență
în rețelistica între containere, aduce conceptul de imagine, astfel putem abstractiza implementarea
microserviciului și să folosim ceva generic indeferent de tehnologia folosită și ne permite
stocare și reutilizare la fel cum am face cu o librărie. De asemenea ne oferă un API pentru
administrarea containerelor indiferent de ce e în interior.

\subsection{Aplicație în container}

Crearea unor multiple containere care funcționează ca un grup poate consuma destul de multe
resurse. Construirea unui simplu proiect în Ktor Server, bazat pe JVM îmi ia aproximativ 2 minute,
astfel dacă am avea o colecție de programe asemănătoare ce se bazează toate pe JVM ar include
rularea a multiplor instanțe de JVM. În acest caz am încerca să optimizăm această structură
și să facem ca să avem un singur contianer cu scopul de a rula întreaga aplicație.

Deși limităm consumul de resurse, anulăm toate eficientizarea ce ar apărea în urma
containerizării. Rulând mai multe aplicații ce ar putea fi rulate independent introducem
probleme în ciclul de viață al aplicației, limităm scalarea și izolarea.

\subsection{Platform as a Service}

Pentru a ușura experiența de lansare există servicii ce ne oferă să ne ia codul, să il
compileze și să creeze un mediu în care să se execute. Acestea abstractizează procesul de deployment
prin crearea unei legături între cod și platformă, pe Heroku aceasta se face prin conectarea
repository-ul sau ca imagine de Docker.

Când nu există probleme ce țin de sistem, nu apare nicio problemă în folosirea unor astfel
de soluții, însă dacă avem nevoie de investigarea la un nivel mai adânc atunci am putea
să vedem cum suntem constranși de aceasta.

Platformele acestea ne pot fi de folos pentru a lansa un serviciu rapid însă putem deveni
constrânși de rigiditatea și de integrarea lor în sisteme complexe. De asemenea, există
mai multe astfel de servicii precum Heroku, Netlify sau Vercel și deși acestea fac același
lucru e posibil ca acestea să ofere facilități diferite.

\subsection{Function as a Service}

O tehnologie ce încearcă se revoluționeze modul de deployment este serverless. Unul din principalele
dificultăți în lansarea codului este provizionarea cu un server pe care acesta rulează,
însă cu tehnologia serverless concentrarea devine asupra scrierii codului iar platforma
pe care o lansăm se va ocupa de rularea și scalarea acestora.

Diferențerea între aceste moduri de deployment este necesitatea de a avea un server
ce lucrează constant în spate pentru a putea prelucra cereri. Cu serverless, codul rulează
doar atunci când e invocat și acest lucru se poate face prin diferite evenimente, fie ele
cereri în browser sau fișiere uploadate în funcție de facilitățile pe care ni le oferă platforma.

Limitările acestei adaptări sunt lipsa de control asupra sistemului, întrucât funcțiile
acestea sunt la interior containere care sunt executate la cerere asta înseamnă ca nu
putem pregăti un mediu pentru acestea. De asemenea acestea nu pot rula la infinit pe toate platformele
existând o limită maximă, acest lucru însemnând că nu putem rula task-uri de lungă durată pe ele.
În domeniul serverless există noțiunea de cold start și warm start. Întrucât funcțiile nu rulează
în continuu, la prima rulare acestea au nevoie de timp de inițializare. O dată inițializate
acestea pot rula mult mai rapid. De asemenea, nu există conceptul de stare, funcțiile
de acest tip nu pot accesa date ale unor rulări anterioare ci doar ce e în instanța curentă.

Cu toate acestea, funcțiile de acest tip ne oferă o flexibilitate în rulare. Dacă în cazurile
precedente, un microserviciu era asociat unei singure instanțe de mașină virtuală sau
de container, la funcții nu e neaparat așa. Putem avea căte o funcție pentru fiecare
serviciu sau un microserviciu format din mai multe microservicii însă pot apărea dificultăți
în partajarea datelor din baza de date.

\section{Infrastructure as a Service}
\subsection{General}
Conceptul de închiriere de echipament nu ar trebui să surprindă pe nimeni. În general, echipamentele
ce ar necesita o investiție semnificativă pentru o folosire extrem de rară ar permite ca împrumutarea
acestora să nu fie o alegere proastă, de exemplu închirierea unei pereche de schiuri în contextul
în care mergem o pentru prima dată sau o dată la o perioadă îndelungată. Compania ce închiriază schiurile
trebuie să se asigure de buna funcționalitatea acestora iar noi plătim prețul pe care l-am negociat.

Însă atunci când vorbim de infrastructură, aceasta este posibil folosită în fiecare moment pe Internet,
rentarea închirierii ar putea provenii strict de la serviciile pe care distribuitorul le-ar oferi, astfel
unul din motivele pentru care ai opta pentru o astfel de opțiune nu este doar pentru componenta fizică
ci și pentru bonusurile care provin cu acestea.

În lumea infrastructurii cloud, serviciile oferite pot fi variate începând de la resurse de calcul sub formă
de platforme de procesare scalabile și actualizabile în orice moment, spațiu de stocare hostat cu servicii suplimentare
precum criptare la rest sau redundanță, rețelistică ușor definită, sisteme de monitorizare și observare până la sisteme 
de Machine Learning deja antrenate alături de alte servicii precum Text-to-Speech sau servicii de traducere în timp real.

Prin implementarea unui serviciu de infrastructură cloud, complexitatea procurării de echipament este abstractizată
prin adăugarea unor resurse noi virtuale, mentenanța la nivel fizic este asigurată de către provider iar în funcție
de serviciul ales chiar mai mult din administrarea acestora, astfel ne putem concentra doar pe dezvoltarea și simplificând
lansarea.

\subsection{Extinderea limitelor}

În comerțul obiectelor cu amănântul, conceptul de cumpărare în masă la un preț mai avantajos sună ușor nedrept
însă furnizorilor să ofere servicii la prețuri foarte accesibile, din acest motiv uneori poate să fie mai 
rentabil să folosim această abordare.

Aceștia pot să creeze ferme de servere pe care să le ofere spre folosire, care deși nu sunt nelimitate, 
ar acoperi suficiente neajunsuri, astfel putem crea aplicații ce beneficiază de această resursă nelimitată în calcul.

În martie 2014, Titanfall a apărut pentru PC, Xbox One și Xbox 360. Consolele noi apărând în 2013, nu este 
neașteptat ca să existe jocuri disponibile pentru ambele, însă în general acestea sunt dezvoltate
ca țintă consolele vechi. În acest caz, dezvoltatorul Respawn Entertainment lucra asupra versiunii de PC și
Xbox One în timp ce un studio extern, Bluepoint, lucra asupra versiunii vechi.

Portarea unui joc ce vizează o platformă mult mai puternică către una mai slabă vine cu dificultățile ei,
mai ales în cazul acestui joc ce este destul de complex (12 jucători iar fiecare dintre ei poate avea 
un titan pe care îl controlează, alături de aceștia existând și entități necontrolabile ce sunt de ambele părți).

În interviul dat pentru Digital Foundry \cite{leadbetter_2014} aceștia menționează dificultățile portării
incluzând reducerea materialelor (asset) și îmbunătățirea modurilor de randare și procesare întrucât 
Xbox 360-ul este mult mai limitat decât Xbox One-ul.

Însă, un lucru special al acestui joc este faptul că este disponibil doar cu o conexiune activă la Internet,
mai mult decât atât, acesta are doar o componentă ce conține mai mulți jucători (fără o poveste), lucru 
ce în absența anumitor cunoștințe ar face să privesc o astfel de abordare ca o lăcomie din partea
distribuitolui.

Chiar și cu o consolă mai puternică, la începutul unei generații dezvoltatorii nu pot să fructifice
toată această putere, astfel aceștia au creat un parteneriat cu cei de la Microsoft ca pentru toate 
versiunile să se folosească servere ce provin din Microsoft Azure pentru a îmbunătății experiența.

Cel mai simplu mod prin îmbunătățirea performanței a fost trecerea la servere dedicate, în general 
jocurile pentru a economisi prețul infrastructurii foloseau o arhitectură de tip peer-to-peer însă 
în cazul Titanfall acesta nu ar fi fost articol, din interviul \cite{leadbetter_2014} am extras:
"The server performance requirements for Titanfall are pretty high, so you can't run peer-to-peer. 
This also means you can't have a many virtual machines per physical server, 
meaning if you've got a popular game, then you're going to need a pretty big investment in 
server hardware.".

În interviul \cite{seppala_2014} se menționează mai multe detalii în acest aspect,
"We have all of this AI and things flying around in the world; that has obviously let us build a 
different game than we would have if we'd have gone with player-hosted", astfel dispozitivele se pot ocupa doar de procesarea informațiilor primite
și să nu se ocupe și de obținerea acestora
"Because Titanfall's advanced AI is handled by the Azure servers, your Xbox's or PC's innards can be used to achieve more detailed graphics 
and the game's silky-smooth frame rate. The Titan bodyguards, 
dropships and legions of AI-controlled combatants are essentially free 
from a processing-power standpoint.".
În același articol se menționează cum serverele disponibile în mai multe regiuni, capacitatea de scalare 
au putut duce proiectul la bun sfârșit.

În articolul precedent \cite{leadbetter_2014} aceste informați despre avantajele cloud-ului se extind,
întrucât pentru a face performanța suficient de bună a trebuit să calculeze anumite elemente direct în cloud.
"Respawn wrote a new PVS system that worked dynamically, likely because that's much better for iteration but it need additional CPU performance we haven't got on 
Xbox 360. We took a different approach and pre-calculated static visibility using a 
big farm of Xbox 360s; these ran multi-hour visibility calculations which reduced the 
amount of stuff the 360 had to draw without requiring runtime computation and allowed us 
to use mostly the same geometry as the Xbox One version.
Same deal for shadows, pre-calculate. Also creating occlusion geometry to reduce 
pixels shaded, tech art to the rescue there getting those built!"

Exemplul precedent acoperă doar resursele de calcul, însă furnizorii de cloud nu se limitează doar la atât,
aceștia oferă servicii mult mai complexe astfel putem să extindem capacitățile unui dispozitiv folosind 
resursele provenite de pe servere la distanță al căror disponibilitate este mai ridicată făță de ce
am putea aduce pe cont propriu.


\subsection{Microsoft Azure}

Principalul furnizor de elemente de cloud este Amazon prin Amazon Web Services, aceștia domină piața întrucât la
început aceștia erau singurii care oferau servicii în care companiile doreau să investească, însă recent,
avantajul lor pe piață este în scădere datorită avansurilor competiției \cite{kumar_2022}, mai ales în lumea
furnizorilor cloud nu este suficient să oferi servicii de hostare întrucât acestea permit o ușoară migrare
către alți competitori în funcție de preț. Din acest motiv, aceștia oferă multe servicii prin care
trecerea devine mai dificilă, sau chiar imposibilă întrucât serviciul este disponibil doar
pe acea platformă, însă în general aceștia oferă servicii asemănătoare, cel puțin la
un nivel de bază.

Am ales să analizez Azure pentru că avem credite ca și studenți. Această platformă
a pornit sub numele de Windows Azure anunțat în 2008. Pe parcursul timpului, acesta
nu a prins atracția publicului larg întrucât era destul de limitată și
tinde să se schimbe foarte des. În perioada aceea Microsoft se baza pe Windows și tehnologii
proprii însă începând cu 2014 Microsoft a început să folosească și să ofere servicii
ce provin din software open-source și își schimbă numele în Microsoft Azure pentru
a reflecta acest lucru.

\subsubsection{Elemente de bază}

Atunci când vrem să folosim un serviciu de cloud, este necesară crearea unui cont,
în general putem folosi un cont Microsoft, fie el personal sau unul administrat de o
organizație.

După crearea unui cont acestuia trebuie să îi se asocieze o sursă de plată, întrucât
acest tip de servicii o să creeze facturi doar pentru cât consumi, astfel administrarea
serviciilor devin mult mai ușoară.

Însă Azure oferă mai multe tipuri de abonamente și modalități de organizare ale acestora
incluzând acces pe bază de rol, destinat companiilor.

Toate elementele create în Azure sunt administrare printr-o interfață REST, Azure
Resource Manager. Acesta poate fi accesat creând cereri HTTP alături de verbele
frecvent utilizate, comunicând prin JSON. Însă pentru a spori eficiența, au fost dezvoltate
SDK-uri de administrare folosite ca librării pentru diferite limbaje de programare
precum Python sau Java, dar și executable-uri făcute pentru a fi rulate direct
dintr-un terminal prin intermediul Azure PowerShell (mai vechi, cu o sintaxă
asemănătoare comenzilor de Windows) sau Azure CLI (mai rapid și cross-platform)

Scopul folosirii serviciilor cloud este eliminarea necesități de organizare
a serverelor, și administrarea acestora printr-un terț ce se oferă să le întrețină.
Un avantaj al acestora este faptul că putem alege locul în care aceste servere
sunt amplasate, nu suntem constranși de locul fizic al acestora, astfel putem definii
Azure Regions.

Acestea reprezintă locul în care cei de la Microsoft au fizic centrele de operare
ale serverelor. În cadrul acestora există mai multe centre de disponibilitate, astfel
aceași zonă prezintă mai multe servere separate ce stau în spatele unui load balancer.
Putem folosi aceste zone pentru a crea replici ca un back-up.

De multe ori avem regiuni pereche, acestea fiind situate într-un loc aproximativ
aproape, de exemplu UK South - UK West, acestea pot fi folosite pentru backup.

Alegerea unei regiuni poate fi ușor dificilă, în primul rând contează locul în
care o amplasăm pentru a scade latența către utilizatori, prezența zonelor de
disponibilitate sau a unei regiuni pereche ca și backup însă datorită modului
cum acestea apar, nu toate caracteristicile sunt disponibile în toate regiunile,
mai mult, toate caracteristiciile oferite de Azure pot să aibă preț diferit
în funcție de regiune, astfel atunci când facem o alegere avem mulți factori
ce ar trebui să îi luăm în considerare.

\subsubsection{Stocare}

În general, tipurile de date cu care interacționăm pot fi de mai multe feluri,
nestructurate, în general obiecte de mari dimensiuni ce reprezintă executabile sau
fișiere de variate tipuri, date semi-structurate, acestea nu au o formă impusă însă
tot au anumite similiratiăți sau moduri de indexare, și date structurate ce reprezintă
date asemănătoare cu cele pe care le-am găsii în bazele de date relaționale.

Modul de interacționare cu stocare se face prin crearea unui cont de administrare
a stocării ce poate fi de mai multe tiprui și are ca scop păstrarea unor anumite setări
de configurare. Acesta poate fi de mai multe tipuri general v2, v1 sau blob
și ne poate limita la tipurile de date pe care le putem folosi. De asemenea,
putem alege modul de stocare al datelor, acestea putând să fie pe hard disk-uri
ce sunt mai lente sau stocare mai scumpă însă mai rapidă prin SSD-uri. În final,
putem selecta modalități de replicare, local (LRS), în care se crează 3 copii în
interiorul unui centru de date din aceași zonă, regional (ZRS), în care se crează
3 copii dar distribuite în mai multe centre de date din aceași zonă, georegional (GRS)
3 copii în mai multe centre din aceași zonă și altele 3 copii într-o regiune secundară pereche
sau georegional doar citire (RA-GRS) în care se primește acces doar pe citire în conturile secundare
din stilul GRS.

Cel mai folosit mod de stocare al datelor este Azure Blob Storage, acesta este un tip
de stocare de date nestructurate, astfel putem sa îl folosim asupra oricarui tip de date,
neavând optimiări specifice. Încărcarea de date prin acest serviciu se poate face grafic
sau prin diferite API-uri. Atunci când folosim acest serviciu, costurile vin din
mai multe categorii: volumul de date, numărul de tranzacții prin interfața REST și
traficul de descărcare. Pentru optimizarea acestor prețuri putem alege între trei categorii,
unul pentru date frecvent accesate, pentru date accesate mai putin frecvente si stocate pentru o perioada
de cel putin 30 de zile și cel de arhivare pentru date accesate foarte rar (mai mult de 180 de zile).

Pe lângă Blob Storage, avem File Storage ce poate să stocheze date sub forma unui drive montabil pe Internet,
Table Service ce reprezintă o bază de date NoSQL de tipul cheie / valoare, fiind o variantă
mai limitată a unei baze de date de acest tip întrucât nu are georeplicare sau scalare, și Queue Service
ce reprezintă o platformă de transmitere de mesaje de tipul Event Bus în arhitectura
microserviciilor.

De asemenea, putem crea Azure Virtual Hard Disks, acestea fiind necesare pentru a
fi montate către mașinile virtuale create.

\subsubsection{Mașini virtuale}

Reprezintă cea mai simplă metodă de a face rost de o platformă de procesare a datelor.
Fiecare mașină virtuală are repartizat o anumită putere de procesare și o anumită
cantiate de RAM, asociat unei unități de stocare și conectat la o anumită rețea.
Toate aceste elemente sunt configurate atunci când creăm mașina virtuală însă pot fi
modificate și dupa ce este lansată.

Putem folosi diferitele metode de acces ale API-ului Azure pentru a modifica aceste date
și putem folosi diferite strategii de migrare precum încărcarea unui .vhd pentru a urca
o mașina virtuală deja existentă.

Modul de monetizare ale mașinilor virtuale este că plătim pentru fiecare secundă
în care aceasta are resurse alocate. În funcție de nevoile noastre putem folosi
diferite tipuri de procesoare în funcțiile de nevoile noastre, fiecare având
optimizările proprii, în același timp putem selecta tipul de stocare de care avem nevoie,
pornind de la HDD-uri standard către SSD-uri ultra rapide, opțiunile mai performante
necesitând o investiție mai mare. Fiecare mașină virtuală necesită cel puțin
o interfață de rețea virtuală și opțional adrese IP publice sau setări de firewall.

Un considerent asupra mașinilor virtuale este modul în care acestea sunt lansate.
Cele mai importante aspecte sunt disponibilitatea și scalarea. Implicit, nu avem
nicio astfel de garanție, din acest moment putem să încercăm fie seturi de disponibilitate
în care mașinile virtuale sunt puse în același centru de date dar în zone fizice diferite
sau o zonă de disponibilitate în care mașinile virtuale sunt puse în regiuni diferite.
Scalarea mașinilor virtuale este mai dificilă, însă există scalare orizontală prin
intermediul Azure Virtual Machine Scale Sets, sunt dificil de implementat necesitând un load balancer, scumpe
deoarece plătim pentru fiecare mașină și necesită configurări și sincronizări manuale.

\subsubsection{Containere}

Întrucât VM-urile au dezavantajele proprii ale consumului de resurse sau al modului
de configurare, putem lansa containere în Azure prin Azure Container Instances.
Acestea sunt configurate aproximativ la fel ca mașinile virtuale, singura diferență
find necesitate specificării numele imaginii containerelui și repository-ul de unde
acestea să provină însoțit de detaliile de autentificare dacă acestea sunt necesare.

Pentru a facilita lucrul cu containere, avem la dispoziție Azure Container Registry
ce ne permite să creăm propriul registru de imagini, putând să adăugăm imaginile proprii
pe el ca după să putem să le consumăm atunci când vrem să lansăm o aplicație sub formă
de container.

Pentru a abstractiza management-ul stării produslui, avem Azure Kubernetes Service
prin care Microsoft abstractizează nodurile de management iar conținutul acestuia
poate fi administrat printr-o interfață similară a lui kubectl, prin az aks.


\subsubsection{Azure App Service}

Reprezentând un PaaS, App Service este o modalitate de a lansare a codului fără
întreținerea explicită a infrastructurii din spatele acesteia, astfel nu e necesar
configurarea unei mașini virtuale și după încărcarea aplicației pe aceasta.

Avantajul renunțării la aceste configurări nu este doar ușurința modului de
lansare ci și posibilitatea unor facilități care nu sunt disponibile mașinilor
virtuale clasice, incluzând replicare și disponibilitate globală, autoscalare dinamică,
integrare nativă în pipeline-uri.

Acest serviciu permite lansare prin intermediul încărcării codului unei aplicații
în mod direct sau prin intermediul unui container. Configurarea se face asemănător
cu o mașină virtuală

Prin intermediul interfeței putem să facem configurări minimale asupra aplicației,
însă acestea sunt limitate și au ca scop mai mult moduri de rulare precum
setări de mediu, porturi sau de SSL.

Monitorizarea unui astfel de serviciu se poate face prin Application Insights
ce se face integrând SDK-ul în codul client, ulterior putem face rost de date
precum timpi de răspuns, performanțe sau telemetrie. Toate aceste date fiind disponibile
în interiorul unui Dashboard central.

\subsubsection{Serverless}

Un concept ce este promovat de către furnizorii de cloud este capacitatea de a lansa
cod fără să ne intereseze locul și modul în care acestea rulează. Deosebirea față
de mașinile virtuale sau App Service este că putem plătii doar cât consumăm, adică
doar timpul în care codul rulează.

Există două tipuri în care putem lansa cod prin această metodă, însă este important
să facem diferențierea de rulare. Tehnologia serverless rulează doar atunci când este rulată,
în spate acestea pot fi sub formă de containere Docker care sunt rulate în momentul
în care e necesar, din acest motiv apare fenomenul de cold-start, atunci când
containerul trebuie să fie inițializat, însemnând că prima rulare a funcției este mai înceată.
În timpurile recente, furnizorii au optimizat acest timp iar rularea unui cold start este mai rară,
acestea fiind păstrare într-o stare în care pot fi rulate aproape instantaneu.

Un alt concept important este faptul că nu avem acces la o stare implicit, fiecare
rulare a funcției vine cu propriile reinițializări, astfel nu putem să accesăm
datele unei rulări anterioare fără să o salvăm conștient într-un anumit loc, precum într-o bază de
date externă.

Azure Functions rulează atunci când sunt invocate, și acesta poate fi configurat
în funcție de necesități precum atunci când sunt invocate prin intermediul unui endpoint HTTP,
la un anumit interval sau la evenimente precum încărcarea unui blob.

Acest tip de execuție, deși abstractizează aplicație, crearea unei funcții
necesită scriere de cod în anumite limbaje compatibile. Acest proces poate fi
simplificat și mai mult folosind Logic Apps, acestea sunt create aproape fără cod
și printr-o interfață grafică. Componentele unui Logic App sunt predefinite prin diferiți conectori
ce integrați împreună vor crea un flux de lucru ce vor fi executate la anumite evenimente.

\subsubsection{Baze de date}

Un alt serviciu pe care îl folosim în crearea unei platforme este o bază de date.
Atunci când avem infrastructură în cloud putem aveam mai multe opțiuni de abordare al
acestei probleme.

Una dintre ele este crearea unei mașini virtuale și în stil clasic crearea serverului
de bază de date și configurarea setărilor de securitate necesare conectării. În acest 
stil putem folosi orice versiune și tip de bază de date însă vom fi forțați să 
ne ocupăm de mașină, astfel va trebui să ne ocupăm de mentenanța software-ului ce
funcționează pe aceasta.

O altă alternativă este folosirea unei opțiuni de tip PaaS pentru baze de date, în 
acest tip toată mentenanța pe care ar trebui să o facem este ascunsă, la fel și 
backup-urile și georeplicarea. Deși pierdem din configurări, avem o metodă de 
folosire a bazei de date mult mai ușoară.

Pentru bazele de date relaționale avem SQL Database ce are suport pentru MySQL, MariaDB
sau Postgres, al cărui tip de monetizare poate fi bazat pe Database Transaction Unit
ce calculează folosința în funcție de CPU, I/O și memorie, vCore ce se face specificând 
un număr de instanțe de mașini virtuale iar numărul acestora va determina consumul și dacă
avem mai multe baze de date avem modelul elastic ce vor cupla mai multe baze de date atunci 
când este înregistrat consumul.

De asemenea, avem opțiunea pentru baze de date nerelaționale prin Azure Cosmos DB
ce are suport pentru multe variante precum MongoDB sau Cassandra ce vor putea avea 
avantajele unei baze de date de acest tip precum georeplicare și flexibilitatea în consistența acestora.

\subsubsection{Monitorizare}


\section{Docker}
\subsection{General}

Docker este una dintre tehnologiile ce a ajutat la revoluția DevOps de a eficientiza
modul de lucru în dezvoltarea software. Acesta a aparut în 2013, și a început să fie
viabil și cunoscut începând cu 2015 însă la momentul de față consider că este o 
unealtă necesară ce nu mai este o noutate și și-a găsit locul în piață.

Containerizarea nu a era o tehnologie noua iar Docker nu este primul și nici singurul
care a contribuit în acest spațiu, totuși aceștia sunt unii dintre cei mai populari 
fapt datorat de ușurința anumitor sisteme noi pe care aceștia l-au adus în favoarea
competiției, aceștia începând trend-ul chiar încă din 2008 prin LXC.

Docker a apărut ca o alternativă la lansarea aplicațiilor în mașini virtuale. Deoarece 
acestea rulau efectiv pe hypervisor și necesitau anumite resurse alocate alături de rularea
întregului sistem ceea ce le făceau greu de configurat. Docker a îmbunătățit acest proces,
aplicațiile care rulează în container sunt mult mai eficiente din punctul de vedere
al consumului de resurse și pot fi lansate mult mai rapid.

Docker permite abstractizarea naturii aplicației, astfel indiferent de ce rulează în ea,
o altă persoană poate să interacționeze cu ea foarte rapid știind doar
un număr limitat de potențiale elemente de configurare. Acest lucru permite separarea 
responsabilitățiilor în procesul de creere. Dezvoltatorii cunosc faptul că aplicația lor va 
rula într-un container, astfel aceștia au responsabilitatea de a containeriza aplicația.
Operatorii cunosc modul de operare cu Docker, astfel aceștia pot administra containerele, 
fără a avea vreo cunoștiințe ce e în ele ci doar de resursele la care acestea ar trebui să 
aibă acces.

Datorită modului în care funcționează, Docker permite modularizarea aplicațiilor, 
astfel fiecare container și-ar rula propria aplicație, astfel poate fi lansat independent, 
fiind susținută o arhitectură bazată pe servicii dar și microserviciile.

\subsection{Componente}

Docker face parte din proiectul Moby, și la mod general reprezintă o aplicație modelată
de arhitectura client-server, în care un executabil Docker va reprezenta un client ce 
comunică cu un server ce reprezintă Docker Engine-ul sau Docker daemon-ul ce are 
responsabilitatea de a rula containerele în funcție de comenziile primite de la clienți.

Nu este necesar ca această aplicație să ruleze asupra aceleiași mașini, astfel putem rula
clientul Docker independent și să folosim Docker Engine-ul pe o altă unitate pe care îl putem 
controla și prin API RESTful.

Componenta de bază reprezintă imaginea. Acestea reprezintă modul în care aplicația noastră
este construită conținând instrucțiiunile de formare. Acestea pornesc de la o imagine de bază 
iar fiecare instrucțiune crează imagini intermediare ce se suprapun sub diferite layere
până se creează complet o imagine.

Registriile reprezintă locul în care imaginile sunt stocate. De aici putem descărca 
imaginile create și să le rulăm. Aceste locuri pot fi publice, precum Docker Hub,
dar și private putând să fie rulate independent. Atunci când le folosim acestea pot 
necesita autentificare, de exemplu pe Docker Hub putem să avem imagini publice 
ce pot fi văzute și accesate de orice utilizator dar și imagini private ale căror 
drepturi de citire sunt restricționate.

Containerele reprezintă o imagine construită ce este rulată, la fel cum un obiect este 
instanțiarea unei clase. Acesta poate avea diferite stări externe precum dacă este creată,
lansată sau distrusă, dar și interne, astfel acestea pot evolua în funcție de modul
cum se interacționează cu ele. Acestea sunt izolate unul de celălalt, astfel chiar dacă
o mașina are mai multe containere acestea nu pot accesa datele unui alt container 
fără să intervenim noi.

Extinderea conceptului de folosire de containere cu Docker reprezintă orchestrarea.
Rularea a unui număr limitat poate fi făcută cu ușurință manual, însă atunci când
avem mai multe aplicații, dependințe (precum baze de date) și alte configurări de rularea 
poate devenii destul de haotic, astfel s-ar pierde multe dintre beneficiile pe care Docker le aducea.
Din acest motiv putem folosii aplicații specifice menite pentru extinderea Docker-ului precum
Docker Compose, Docker Swarm sau Kubernetes pentru a administra sisteme mai complexe.

\subsection{Utilizare}

Din punctul de vedere al unui dezvoltator, folosirea Docker-ului se rezumă
doar la câteva instrucțiuni de bază ce pot fi cuprinse pe o listă foarte 
\href{https://dockerlabs.collabnix.com/docker/cheatsheet/}{scurtă}, din acest 
motiv o prezentare în detaliu al fiecărei comenzi nu ar fi eficientă.

Însă, pot sublinia câteva utilizări generale după instalare.

\begin{itemize}
    \item Pornirea oricărei aplicații ce are o imagine publicată pe Docker Hub,
precum Ubuntu, \verb|docker container run ubuntu| va încerca să lanseze un container 
ce rulează din imaginea ubuntu. Dacă nu găsește imaginea ubuntu acesta va încercă să
găsească imaginea ubuntu cea mai recentă (latest) de pe repository-ul implicit (Docker Hub),
dacă o va găsi acesta o descarcă și o să ruleze
    \item Atunci când o imagine este rulată, acesta este executată cât timp are un proces în execuție,
din acest motiv, containerul anterior de Ubuntu se va închide imediat ce va porni. Putem accesa 
și prelungii viața containerului prin preluarea STDIN-ului și atașarea unui pseudo tty folosind
\verb|docker container run -it ubuntu bash|, acesta va rula alt container, și ne vom găsii
în interiorul unui terminal cu root@containerid.
    \item Când lansăm o imagine, implicit acesta primește o rețea de tipul bridge, astfel 
aceasta poate fi conectată la rețeaua computerului principal. Putem schimba acest aspect
prin directivele din -network ale comenzi \verb|docker container run|.
    \item Containere diferite nu pot implicit să folosească datele altui container. Din acest motiv,
Docker ne permite să configurăm Volumes, acestea au rolul de a salva anumite date și să fie accesibile
sub diferite forme. Putem lăsa Docker să le administreze sau putem crea bind-mounts care pot să facă 
ca fișierele sau un director de pe mașina gazdă să fie sincronizate cu un director din interiorul mașinii.
Acestea se fac folosind directivele -v ale comenzi \verb|docker container run|.
    \item În general am vrea ca aplicația din interiorul containerelui să fie accesibilă din exterior. Acest lucru
se poate face sub diferite forme, însă în general vrem să configurăm un port de acces. Astfel aplicația
din interior comunică pe un anumit port (de exemplu 3000), însă Docker nu o să deschidă implicit aceste 
porturi. Pentru a face asta, trebuie să folosim directiva -p care va asocia port-ul mașinii gazdă cu port-ul
din interiorul containerului. De exemplu \verb|docker container run -p 3001:3000|, port-ul 3001 pe gazdă va accesa aplicația din container
care rulează pe portul 3000.
    \item Containerele sunt puse într-o rețea generală de către Docker, însă putem crea rețele proprii 
și să agregăm mai multe containere în el, astfel nu este nevoie să oferim acces către exterior tuturor 
serviciilor ci doar celor strict necesare.
    \item Construirea unei imagini se face, în general, folosind un Dockerfile, acesta pornește
de la o imagine de bază asupra căreia se adaugă diferite configurări precum comenzi sau specificare
de argumente de build sau environment. Comanda este \verb|docker image build .| ce va 
construii imaginea din Dockerfile din directorul curent. Putem taga imaginea folosind directiva 
\verb|-t|.
    \item Încărcarea unei imagini într-un repository se poate face folosind \verb|docker image push|.
\end{itemize}

\subsection{Simpla orchestrare}

O metodă ușoară de a crea un sistem format din mai multe containere este de a folosi
Docker Compose. Întrucât Docker nu a fost făcut pentru a administra mai multe containere
în același timp și în același timp acesta susține ca unui singur serviciu să îi corespundă
un container, putem avea dificultăți atunci când am vrea să avem un sistem mai complex,
de exemplu un server web și baza de dată asociată acestuia.

Docker Compose este o metodă de orchestrare, astfel acesta crează metode de rulare de containere,
astfel dacă vrem să avem o aplicație în Compose, acesta trebuie să fie deja 
sub forma unei imagini.

Pentru a începe, avem nevoie de un fișier .yml. În acesta putem definii diferite
servicii ce reprezintă defapt containerele pe care vrem să le rulăm. Pentru fiecare
acest serviciu putem definii în acest fișier toate proprietățile pe care le-am putea
definii atunci când am rula \verb|docker container run|, precum volume, rețele, porturi.
De asemenea putem definii (pornind din imagine) health check-uri, comenzi care ar fi rulate
în interiorul containerelui și diferite proprietăți care ar definii starea unui serviciu, astfel
în funcțiile de setările făcute, dacă un serviciu de care depindem nu este disponibil, am putea
închide toate serviciile care depind de acesta, sau să încercăm resetarea serverului afectat.

Docker Compose ne permite să avem o imagine mai per ansamblu asupra sistemului pe care vrem să îl creăm.
Însă acesta are limitările proprii ce îi oferă imaginea de o metodă simplă, una dintre acestea
fiind faptul că nu avem asistență pentru replicare sau metode de actualizare ale versiunii unei imagini ce
pot fi găsite în versiuni mai avansate ale sistemului.

O metodă de a face rost de replicare într-un mod simplu este folosirea unui serviciu
de service discovery precum Hashicorp Consul, astfel putem avea un cluster pe care
să rulăm diferite containere replicate, însă nu am cunoaște exact locația acestora. În același 
stil, locația acestora ar putea să se schimbe etc. O astfel de unealtă folosește metode 
asemănătoare cu un DNS pentru a crea comunicare în cadrul unui cluster, astfel 
am putea să avem un load balancer și prin intermediul Consul să se ruteze 
către replica cea mai avantajoasă. Specific acestui servicii putem să beneficiem de 
mutual TLS pentru a securiza rețeaua.

O altă metodă de orchestrare simpla reprezintă Docker Swarm, aceasta vine la pachet 
cu Docker Engine, având un mod special pentru acest tip de operații. În acest mod, 
un Docker Engine se poate alătura(sau crea) unui cluster de alte noduri care rulează
în acest mod. Fiecare nod va fi fie Manager sau Worker, având minim un master(inițial,
primul nod este nodul manager). Nodurile manager au rolul de a distribuii task-urile primite
către nodurile worker, însă acestea pot efectua task-uri la fel ca un nod obișnuit,
însă aceasta nu este o prioritate pentru ele sau poate fi dezactivată în totalitate.
Un avantaj al acestui serviciu este posibilitatea de a rula replici ale aceluiași serviciu,
astfel putem ca un singur container (serviciu) să ruleze simultan în mai multe locuri, 
astfel dacă un loc este afectat le putem accesa pe celelalte, funcționând ca un load balancer.
O limitare al acestui sistem este faptul că această scalare nu este automată și că
sunt metode limitate de monitorizare implicit însă este extrem de ușor de configurat și de rulat,
astfel pentru unele sisteme reprezintă o alternativă suficient de bună a unor soluții mai complicate.

\section{Docker Swarm}
\section{Kubernetes}
\section{Red Hat OpenShift}
\section{HashiCorp Nomad}
\chapter{Îmbunătățirea metodelor de dezvoltare software}
\section{Introducere}
În dezvoltarea aplicațiilor au fost realizate diferite optimizări în modul în care acestea
sunt create pentru a mării productivitatea. Acestea pot avea diferite forme,
precum sistem automat de testare, build, deployment sau orice altă cerință ce
în general e făcută manual ce se găsește o metodă automatizată pentru a fi efectuată.

În cazul aplicațiilor cu o arhitectură bazată pe microservicii sunt esențiale
pentru a avea un progres acceptabil în dezvoltare. Cu cât sunt introduse mai multe
microservicii, cu atât există mai multe elemente ce pot să se strice atunci când este
efectuată o modificare. Din acest motiv, este necesar să aducem îmbunătățiri în
modul cum dezvoltarea este realizată și să detectăm problemele într-un timp scurt.

Cea mai populară dintre utilizăriile cunoștințelor acumulate prin învățarea 
programării probabil este dezvoltarea de aplicații software în diferite forme, 
fie în cadrul unei interprinderi sau individual. Pentru a avea orice fel de succes,
aceasta trebuie să aducă valoare unui public iar această necesitate este constrânsă
de diferite aspecte manageriale și de marketing, exprimând în mod simplu, trebuie
ca produsul să răspundă la nevoile clienților în timp util.

Crearea de produse software pentru consumatori sau interprinderi, este o practică ce 
poată să fie considerată suficient de nouă, atribuită perioadei în care computerele 
au putut să fie accesibile și pentru aceștia iar ei vedea valoare în achiziționare 
a acestuia.

Astfel, la fel ca alte aspecte din industria IT, multe concepte au fost preluate
din practici deja existente, chiar dacă cumpărarea de aplicații software sau utilizarea
acestuia este recent, oamenii au fost nevoiți să negocieze alimente sau obiecte
pentru diferite întrebuințări de foarte mult timp. Însă datorită schimbărilor
ce au dus la dezvoltare tehnologică dar și a formei în care aceste aplicații sunt consumate
s-a observat o constantă nevoie de îmbunătățire ale livrării pentru a 
răspunde la nevoile clienților în timp util.

Însă concentrându-ne doar pe produs și pe modalitatea creării acestuia poate să
nu fie benefic. Produsele software au început să devină complexe, nu mai sunt
de independente și au diferite dependințe, ceea ce determină necesitatea apariției
comunicării între echipe diferite ce se ocupă de părți diferite ce susțin
buna funcționare a proiectului.

Din acest motiv, de-a lungul timpului au apărut diferite tendințe dar și idei 
ce au ca principal obiectiv îmbunătățirea modului în care interacționăm cu procesele
și componentele ce duc la formarea produsului final. 
\section{Metodologii de dezvoltare software}
Eșecurile în proiectele software sunt scumpe, întrucât acestea necesită investiții
semnificative. De aceea optimizarea modului în care acestea sunt create este esențial,
din acest motiv s-a ajuns la coneperea și adaptarea a diferite metodologii
și framework-uri ce au ca scop oferirea unei baze de inspirație
a unui mod de organizare ce poate fi considerat că ar crește șansele
ca proiectul să fie un succes.

\subsection{Waterfall}

Orice obiect destinat consumului, de exemplu diferite unelte sau alimente,
au un proces standard de dezvoltare care cuprinde într-un mod foarte simplist
o linie de ansamblare în care se introduc materialele necesare producerii
iar multe persoane contribuie pentru ca în final să se obțină un singur produs final.

Aplicațiile software sunt diferite întrucât acestea funcționează mai asemănător
cu un serviciu, acesta este creat o singură dată iar apoi este distribuit
către consumatori. Astfel, munca efectivă este creată o singură dată pentru
ansamblarea acestuia.

Însă de-a lungul timpului, modalitatea de distribuție a fost cea care a evoluat,
dacă la început livrarea produselor era foarte asemănătoare cu livrarea
produselor obișnuite, în care exista un produs declarat final și care
era livrat sub formă imutabilă, de exemplu pe CD-uri, acum livrarea se
face doar încărcând produsul într-un anumit loc, sau lansând direct
într-o platformă cloud.

Observăm faptul că o dată ce lansăm produsul, modificarea unor eventuale
probleme este complexă, necesită retrimiterea prin formatul inițial și eventual
chemarea înapoi a lucrurilor deja livrate, acest proces este complex și costisitor.

Din acest motiv este necesar ca produsul dorit să fie descris clar
pentru a elimina eventualele modificări ce ar putea apărea și trebuie să fie
bine testat întrucât corectarea unor greșeli dupa livrare devine imposibilă,
astfel metodologia Waterfall este liniară, are un stagiu de start și unul
de final și se axează pe documentație ce descrie produsul și pe validarea acestora.

Etapele metodologiei Waterfall:

\begin{enumerate}
    \item Crearea specificațiilor, include comunicarea între client și echipă
          a condițiilor ce proiectul ar trebui să le îndeplinească. Acestea trebuie să
          fie foarte clare și să nu lase aspecte ce pot fi interpretate întrucât
          acestea teoretic nu pot fi schimbate în perioada dezvoltării
    \item Crearea arhitecturii, include interpretarea cerințelor clientului
          în vederea creării design-ului proiectului și alcătuirea specificațiilor tehnice ale
          produsului
    \item Implementare
    \item Validare
    \item Mentenanță, după validare proiectul este trimis către client iar acesta
          este considerat final, însă clientul poate revenii pentru a aduce îmbunătățiri
          produsului.
\end{enumerate}

\subsection{Agile}

Metodologia Waterfall este rigidă și nu permite modificări în timpul dezvoltării,
însă acest aspect nu este compatibil cu dezvoltarea software întrucât cerințele
clienților în general nu sunt clare sau se schimbă. În același timp unele
caracteristici trebuie să fie livrate la un anumit timp pentru a avea efectul dorit,
astfel dacă perioada este depășită nu pot să aibă același impact.

Însă în același timp, nevoia de livrare a unui produs complet a scăzut,
întrucât costul de livrare nu mai este imutabil a scăzut necesitatea
verificărilor intense și alcătuirea documentației din Waterfall întrucât
acestea se pot schimba.

Agile se bazează pe crearea unui Minimum Viable Product (MVP) cât mai rapid,
astfel clientul poate să aprecieze modul cum arată produsului.

Pentru a capta aceste schimbări, dezvoltarea în agile se face prin intermediul
unui set de principii și valori ce se axează pe crearea unui mediu în care schimbările
sunt acceptate în orice moment iar părerea clientului este importantă, de aceea
livrarea către acesta trebuie să fie cât mai desă iar produsul trebuie să fie
mereu într-o stare în care poate să fie livrat.

Agile scrum este un framework ce aduce adăugări precum ședințele de planificare,
retrospectivă, ședințe zilnicie în care echipa comunică la ce lucrează și
eventualele probleme pe care le au dar și diferite roluri din cadrul echipelor.

\subsection{DevOps}

Deși nu ar trebui să fie încadrată ca o metodologie, DevOps este în primul
rând o mișcare culturală ce încearcă să rezolve unele probleme de natură
socială ce apar în dezvoltarea aplicațiilor software.

Poate fi considerată un super-set al Agile întrucât majoritatea principiilor
declarate se pot aplica și atunci când încercăm să ne orientăm spre adaptarea
la o astfel de acțiune.

Întrucât nu există o definiție clară, poate fi deficil de estimat dacă
adoptăm DevOps corect, în general putem considera că dacă folosem anumite
practici precum CI/CD sau încercarea optimizării modurilor de lansare am putea
să considerăm că este suficient, însă este incomplet întrucât DevOps
promovează mai mult cultul organizațional și armonizarea relațiilor dintre
persoane care duc la îmbunătățirea perfomanțelor.

\section{Aspecte organizatorice și culturale}
Primul punct pe care DevOps se concentreaza sunt aspectele culturale ale companiei.
Însă acest lucru nu este nou, management-ul poate fi considerată o ștință
ce se ocupă cu descoperirea aspectelor desfășurate într-o companie ce o cauzează
să aibă succes, iar o mare parte din acestea o reprezintă relațiile pe care
angajații o au la locul de muncă.

DevOps a fost elaborat pentru a sporii randamentul pe care echipele o au atunci
când interacționează cu produsul. Însă orice aplicație probabil se va conecta
la o bază de date prin intermediul unui serviciu web, tipica arhitectură pe trei straturi.
Din acest motiv, multe echipe cu roluri diferite trebuie să interacționeze,
dezvoltatorii, testerii, oamenii ce se ocupă de lansări (operatori), persoanele
ce administrează infrastructură, alături de aceștia sunt diferite persoane din
lumea de afaceri, manageri, deținătorii de produs, agenții de vânzării, echipele de suport
și alti stakeholders.

Atunci când avem multe echipe independente ce lucrează pe un singur produs
poate devenii dificil să lansăm caracteristici noi prin metode clasice
dacă ignorăm modul în care tratăm produsul.

Istoria cuvântului DevOps, în mod simplist am putea spune că provine de
la armonizarea dintre Dezvoltatori și Operatori. Acest grup de persoane
pot fi considerate că sunt „dușmani” ce sunt forțați să interacționeze
pentru a crea produsul final. Relația această este creată de modul în care
cele două roluri sunt promovate, Dezvoltatorii vor fi mereu celebrați pentru
crerea de caracteristici, rezolvarea de probleme și sporirea potențialui
valori pe care produsul creat ar putea să o aducă, din acest motiv
aceștia vor să lanseze în piață cât mai des. Operatorii sunt promovați
pentru stabilitatea produsului, astfel clienții trebuie să poată să își acceseze
serviciul în orice moment, astfel este pusa foarte mult presiune pe aceștia
ca serverele să fie operabile, iar lansările sunt punctele cele mai critice
întruât orchestrarea acestora poate să fie dificilă și poate să cauzeze
probleme pentru consumatorii.

O astfel de relație apar și între dezvoltatori și echipa de testare,
vizibil în modul Waterfall unde dezvoltatorii trebuie să se întoarcă
să lucreze la un proiect pe care ei îl considerau terminat, sau între
dezvoltatori și echipele de securitate ce testează vulnerabilitățile.
Întrucât acest tip de verificare este făcut abia la final, reînceperea
dezvoltării pentru a fixa anumite probleme poate să fie costisitoare.

Armonizarea unei astfel de relații este cea care această mișcare dorește
să o rezolve, și nu este singura. Studiile de management dar și intuiția
ne pot arăta că mediul în care lucrăm ne ajută să performăm. Dezvoltarea
aplicațiilor reprezintă o metodă de inginerie ce este extrem de taxantă asupra
persoanelor implicate, din acest motiv crearea unui mediu de lucru ideal
este importantă iar aceasta poate pornii de la cultul organizațional.

Însă fixarea unor astfel de probleme necesită ca mediul de lucru să
permită schimbări, fie ca managerii să accepte o perioadă de ineficiență
iar persoanele implicate să fie comunicativi și să accepte reunțarea
la unele obiceiuri vechi ce funcționează însă care pot să fie considerate ineficiente.

O astfel de metodă poate fi împuternicirea echipelor prin creare unor echipe
funcționale cross funcționale ce le permite participanților să facă mai multe
elemente din cadrul produsului, astfel cunoaștem mai multe părți ale acestuia.

Dezvoltarea software pornește de la lucrurile comune pe care le știam, astfel
putem să aplicăm strategiile pe care companiile de dezvoltare le aplică asupra
construiri obiectelor dificle precum mașinile.

O metodă ce a revoluționat dezvoltarea în Japonia, a fost producția de tip lean
apărută la Toyota ce se axează pe maximizarea productivității și minimizarea
pierderilor, adică a tuturor aspectelor ce nu generează valoare
pentru client. Pe scurt, aceasta se bazează pe un proces de
dezvoltare continuă pe baza unui set de principii, identificarea modalităților
de apariție a valori, identificarea circuitului necesar pentru crearea valorii
de la obținerea materiilor prime până la produsul final, crearea unui
flux pentru a minimiza timpii morți și a unui sistem de creare a strictului
necesar pentru a nu consuma timp inutil pe elemente ce nu aduc valoare
și perfectizarea întregului proces în orice moment pentru a rămâne competitivi.

Sunt multe aspecte pe care le putem desprinde din acest mod de dezvoltare
și să îl aplicăm în tehnologie, însă mi se pare că cel mai important dintre
acestea reprezintă crearea acestui circuit al valorii. Atunci când am vorbit despre
această dezvoltare de software am vorbit despre aceste echipe diferite
ce lucrează toate pentru a aduce valoare printr-un singur produs.

Există multe situații în care rezolvarea unor probleme necesită atingerea
unor puncte diferite din cadrul sistemului, astfel o abordare DevOps dorește
să împuternicească persoanele implicate, dându-le posibilitatea de a cunoaște
și alte sisteme pe lângă cele asupra celor care lucrează. De exemplu, fiind
dezvoltator a unei aplicații software ce încearcă să rezolve o problemă
de natură de server, posibilitatea de a vizualiza codul de server
pentru a identifica posibile probleme ar putea să ajute. Astfel captarea
circuitul ce crează valoare pentru noi ar putea să însemne să avem acces
asupra tuturor elementelor pe care le folosim, de la infrastructură, la servere și baze
de date pentru a nu crea timpi morți așteptând alte persoane să rezolve probleme
ce ne impactează în mod direct.

Îmbunătățirea performanței nu se rezumă doar asupra lucrurilor care
generează profit, dar și menținerea echipelor „în formă” din punct de vedere
mintal. Dezvoltarea de software include probleme ce sunt taxante din punct de vedere
cognitiv, din acest motiv crearea unui mediu relaxant poate ajuta în randamentul personalului.
Însă în același stil putem dorii să sporim creativitatea angajaților. Multe
din diferite idei pornesc de la echipele de inginerie ce se confruntă cu anumite
probleme și vor să creeze soluții pentru a le rezolva, de exemplu la Gmail a pornit de la
un angajat frustrat cu soluția de mail pe care trebuia să o folosească și a dezvoltat
în timpul propriu acea problemă. Unele companii își asignează timp, o zi pe lună,
o săptămână la câteva săptămâni în care angajații lucrează la aceste proeicete
secundare care probabil nu vor genera profit însă care îi ajute să se dezvolte.

Un alt aspect al cultului organizațional ce se poate schimba reprezintă modul
în care aceasta tratează comunitatea, acest lucru referindu-mă la soluțiile open source
pentru diferite produse dar și interacțiunea cu alte companii.
În mediul de afaceri, menținerea avantajului competitiv reprezintă una dintre
prioritățile pe care se axează, pentru a ne diferenția de competitori. Însă
acest lucru nu ar trebui să se axeze și asupra problemelor pe care le întâlnim,
dezvoltatorii ar trebui să fie încurajați să își exprime problemele pentru a
evita alte persoane de a se bloca asupra lor. În același mod folosirea de
pachete open source duce la optimizarea timpilor de dezvoltare dar și la
calitatea produselor.

O altă practică benefică reprezintă motivarea angajaților în mod corespunzător.
Un angajat fericit va avea randament bun. Managerii consideră că partea financiară,
salariul și diferite bonusuri, sunt suficiente pentru a păstra angajații însă
nu este suficient. Un salariu corespunzător nu reprezintă o așteptare, ci o necesitate,
astfel aceștia ar trebui să fie plătiți minim prețul pieții pentru poziția lor.
Dar și ascultați la diferite nivele, fie prin diferite chestionare sau direct,
aceștia pot fi motivați în diferite forme în funcție de dorințele acestora
iar o coordonare bună va folosi dorințele angajaților în atingerea obiectivelor companiei,
însă acesta nu reprezintă neapărat un aspect legat de DevOps ci unul de bună practică în management.

Monitorizarea este un alt aspect important în DevOps, însă nu ar trebui să fie
o surpriză. O companie este evaluată asupra diferiților indicatori de performanță,
iar aceasta se va extinde și la echipele din cadrul acestora. Însă indicatorii
de performanță la nivel de echipă pot fi diferiți, întrucât aceștia au scopul
de a îmbunătății perfomanța echipelor, iar colectarea acestora poate fi făcută în mod diferit
în funcție de procesele firmei. Însă aceste date nu ar trebui să fie folosite
pentru a evalua angajații ci strict pentru a îmbunătății perfomanța echipei.
Atunci când vorbim de performanța unei singure persoane, aceasta nu se traduce
în performanța echipei, iar în general putem să ajungem la crearea de procese
cu scopul perfecționării.

Un aspect important al culturii reprezintă modul de gândire asupra eșecului.
În general, companiile ar încerca să ascundă și să prevină orice problemă cu
produsul pe care îl crează, însă este ineficient să gândim în acest mod, din același
mod în care nu ar trebui să eficientizăm un proces pe care nu știm dacă are nevoie să
fie eficientizat. Astfel o metodă mai bună de combatere a problemelor asociate
produsului este să creăm procese rapide de redresare a acestora prin posibilități
de întoarcere la versiuni anterioare sau de aplicare a unor hotfix-uri cât mai rapid, astfel scăzând
timpul mediu de reparare (MTTR) dar încercand să minimizăm timpul mediu dintre eșecuri (MTBF)
și timpul mediu de eșec (MTTF), de asta e necesar să colectăm fișiere de jurnalizare
și metrici și implemetarea de sisteme ce au ca scop analizarea acestor date
și să încerce să prezică posibilele eșecuri. De asemenea este important să
implementăm sisteme ce minimizează impacul creat de apariția unui erori ce au ca scop
minimizarea posibilelor pierderi din apariția lor, prin diferite
metode precum eșuare rapidă prin care componentele din jur detectează eroare unui sistem
de care depind și își schimbă modul în care se comportă dar și implementarea
unui sistem eficient de avertizare ce poate să ilustreze situația reală a sistemului
și gravitatea acestuia.

Eșecul este inevitabil, iar modul în care îl tratăm este important. Un alt aspect important
în cultura organizatorică este păstrarea unei atitudini în care nu se atribuie vina către
o anumită persoană („blameless”), oamenii fac greșeli și nu le fac în mod intenționat.
Apariția unui incident trebuie să fie urmată de o ședintă de revizie a modului
cum a decurs. Aceasta poate include mai multe detalii precum crearea unui cronologii
a evenimentelor ce s-au desfășurat. Însă aceasta nu va ajuta la descoperirea unei cauze
principale, ce este o practică ce ar trebui să fie evitată. J. Paul Reed mentionează
că găsirea unei astfel de cauze reprezintă doar locul în care te oprești din a investiga mai departe.
Aceste ședințe ar trebui să fie urmate după rezolvarea incidentului și să aibă ca scop
doar îmbunătățirea procesului pe viitor și nu ca motivație negativă. Personalul
ce este amenințat nu va încerca să facă schimbări, astfel nu vor avea oportunitatea de a
crește și de a se perfecționa.

\section{Aspecte tehnice}
Un aspect pe care se concentrează majoritatea persoanelor atunci când vorbesc
despre DevOps reprezintă uneltele sau aspectele practice pe care urmarea
acestei mișcări o susține, însă nu este suficient mai ales că folosirea
unor unelte nu garantează o creștere a eficienței.

Însă un lucru este cert, pentru a crește eficiența din cadrul unei echipe este
necesar să automatizăm cât mai multe aspecte, de aceea microserviciile nu pot fi
folosite eficient fără a avea o bază formată ce automatizează aspecte precum
testarea și lansările.

Pornind de la nivelul cel mai jos, folosirea unei metode de versionare a codului
reprezintă un standard, aceasta se face prin intermediul Git, iar hostarea repository-urilor
poate să fie făcută pe o platformă precum GitHub, GitLab sau BitBucket. Toate acestea
ar trebui să fie accesibile de orice persoană ce are ca responsabilitate un produs.
Aceasta ne poate ajuta să cunoaștem alte produse și să aducem modificări, chiar și din
exterior atunci când e nevoie prin intermediul pull request-urilor.

Pentru a elimina probleme ce apar datorită compilării pe un dispozitiv local,
putem folosii Maven pentru a automatiza procesul de build, astfel putem să fim siguri
că programul este lipsit de dependințe locale.

Stocarea și versionarea produselor rezultate după build, poate fi făcută
cu un program asemănător cu Nexus Artifact Management, acesta poate fi folosit
pentru fișiere jar, war sau pachete npm și oferă suport pentru controlarea
accesului. Acesta ne permite să reutilizăm rezultatele pasului de build în mai multe locuri.

Un alt aspect pe care putem să îl automatizăm conturează integrarea continuă, o astfel
de unealtă reprezintă Jenkins ce este folosit pentru a sincroniza branch-urile,
rularea de teste automate asupra lor și eventual construirea artefactului și trimiterea
acestuia spre stocare.

Automatizarea infrastructurii și menținerea acesteia sub formă de cod reprezintă
un alt proces ce poate să fie automatizat. Menținerea sub această formă ne aduce
avantaje precum notificarea tuturor personalor implicate prin intermediul pull request-urilor
dar și crearea unui istoric ce poate fi folosit în cazul unui audit. Pentru a automatiza
modul în care procurăm infrastructură putem folosii Terraform, o unealtă folosită
pentru a crea resurse în cloud prin intermediul codului oferind starea pe
care vrem să o avem. Acesta poate fi folosit pentru majoritatea platformelor cloud dar și
peste echipament on-premise.

O dată ce echipamentul este procurat, o altă dificultate reprezintă configurarea acestora. În trecut
acesta reprezenta un lucru manual extrem de repetitiv, însă atunci când avem de a face
cu un astfel de proces pot exista probleme ce apar din erori umane, astfel
automatizarea ar ajuta prin evitarea acestora. O altă unealtă pe care putem să o folosim
în acest scop reprezintă Ansible, ce ne ajută în configurarea automată a infrastructurii
prin intermediul codului.

Atunci când trebuie să folosim părți de aplicație ce nu au un artefcat clar (de
exemplu un server în Node.JS) putem folosi containere Docker pentru a le stoca
sub forma unui container, astfel acesta poate să fie portat dintr-un mediu în altul și
să fie refolosit la mai multe stagii.

Pentru a orchestra modul în care containerele rulează putem folosii Kubernetes, acesta
ne permite să administrăm cluster-ul și avem facilitatăți precum DNS în interiorul acestuia,
posibilități de recuperare și orchestrare a lansărilor.

Monitorizarea reprezintă un alt aspect important pentru a estima starea sistemului,
de aceea putem folosi ceva asemnănător cu stack-ul ELK. Elasticsearch reprezintă
o bază de date cu posibilitate de căutare pentru a stoca elementele de jurnalizare.
Logstash oferă o metodă de procesare a elementelor de jurnalizare pe care le
preia din diferite surse și pe care le trimite în Elastic. Kibana reprezintă o platformă de vizualizare a
datelor stocate în Elastic.

Pentru a automatiza testarea putem folosii framework-uri precum Selenium, astfel
aceste teste vor rula de fiecare dată când se creează un build. În același stil
putem folosi unelte de analizare a seucrității prin integrarea acestora în interiorul
ciclului.

Testarea se face la diferite nivele, astfel testarea automată poate să aducă
un sens de securitate pentru integrarea schimbărilor, acestea fiind vizibile
pe pull request. În același stil, pentru testarea manuală putem să creăm medii efemere
ce sunt create strict pentru a testa acel branch. Pentru a face asta putem folosii Vagrant.

Atunci când avem multe sisteme, va exista necesitatea administrării certificatelor,
parolelor și cheilor de API externe, acestea trebuie să aibă o anumită perioadă de timp
cât sunt valabile iar în cazul în care sunt compromise trebuie să aibă posibilitatea
să fie schimbate și să se reflecte acestu lucru în sistemele deja existente. Pentru a rezolva
această problemă putem folosii HashiCorp Vault ce oferă aceste funcționalități.

Chiar dacă sunt mai puțin vizibile, sistemele de comunicare sunt altă parte importantă a
DevOps-ului, fie că acestea reprezintă un sistem de Mail sau o platformă de
comunicare în timp real precum Slack, acestea trebuie să îmbunătățească posibilitatea
de sincornizare a muncii angajaților dar și de crearea unei căi de alertare în cazul
apariției erorilor. În același timp, acestea trebuie să fie folosite pentru a păstra
la curent persoanele care nu pot participa la ședințe sau lucrează de la alte sedii sau de
acasă.

O bună parte din automatizări se referă la infrastructură și la configurarea
acesteia, din acest motiv este util dacă folosim servicii administrate
de către providerii cloud prin platformele acestora. Acestea oferă multe alte
servicii pe lângă sistemele clasice și pot oferii servicii de notificare și monitorizare complexe
ce ușurează ciclul de dezvoltare.

\section{Aspecte aplicabile}
Chiar și în timp ce citeam despre DevOps toate elementele din jurul acestuia
deveneau din ce în ce mai neclare. Aceasta poate să fie și din faptul că
discută despre elemente pe care eu nu am putut să le experimentez personal, mai ales
cele legate de organizație sau de unelte al căror întrebuințare depășesc nevoile
proiectelor asupra cărora vreau să lucrez.

Din acest motiv, poate să fie ușor neclar ce reprezintă DevOps, și asta poate să
fie văzut și în resursele disponibile, întrucât nu există o anumită formalitate
asupra acestui termen.

Un lucru este cert, DevOps este o mișcare ce dorește să mobilizeze personalul
de a îmbunătății modul cum interacționează cu procesul de dezvoltare prin diferite forme,
îmbunătățirea comunicării, creșterea nivelului de automatizare și delegarea responsabilităților
către diferite servicii.

Astfel, atunci când vrem să adoptăm acest concept este suficient doar să îmbunătățim treptat
lucrurile la care lucrăm și care ne influențează, ignorând mentalități precum
păstrarea modurilor antice dar care funcționează pentru a crește randamentul procesului.


\chapter{Aplicație}
\section{Scopul aplicației}

Chiar dacă o aplicație de tipul acesta nu ilustrează capabilitățile unui sistem
bazat pe microservicii, în urma documentației putem spune că din punct de
vedere arhitectural, folosirea unei acestui tip nu ar aduce
aproape niciun avantaj întrucât există un singur dezvoltator deci segmentarea
ar duce complexitate inutilă iar populația
care va folosi aplicația va fi aproape zero, astfel nu avem nevoie de scalarea
independentă a unor componente. În același timp, deoarece suntem constrânși de resurse
cea mai bună metodă este să delegăm cât mai multă administrare către cloud,
astfel să putem să ne concentrăm cât mai mult pe dezvoltare.

Cu toate acestea, aplicația pe care vreau să o dezvolt are ca scop rezolvarea
unei probleme personale, crearea unei aplicații de administrat task-uri cu diferite
capabilități precum jurnalizarea. În mod mai simplu, vreau să fac o clonă
de Atlassian Jira însă mai simplistă, dedicată oamenilor obișnuiți ce nu au nevoie
de complexitatea unui program de administrare de proiecte.

Ceea ce m-a atras la această idee reprezintă ideea de „Log work”, prima oară
când am interacționat cu această mecanică mi s-a părut interesantă întrucât
se crează un istoric al tuturor activităților pe care le fac, putând
să îmi creez o imagine atunci când mă uit în retrospectivă și să planific mai
ușor în viitor. Alături de aceasta, faptul că am task-uri împărțite sub diferite forme mă ajută
să îmi organizez ziua mult mai eficient.

Jurnalizarea este individuală iar modul în care fiecare vrea să își organizeze
task-urile reprezintă o preferință personală, din acest motiv este dificil să
găsești o aplicație care să se potrivească, din acest motiv nu prea am folosit aplicații
deja create întrucât mi s-a părut că sunt ușor influențat de modul cum ar trebui să
îmi administrez cerințele.

Prin această aplicație vreau să îmi organizez ziua mai ușor, să creez o soluție
de aministrare a lucrurilor care vreau să le fac zilnic ce poate să le aranjeze
automat într-un mod inteligent pentru a-mi oferi o perspectivă asupra zilei.

\section{Aplicatii asemanatoare}

\subsection{Jira}

Jira este una dintre aplicațiile pe care trebuie să le folosesc zilnic
la muncă și din experiența proprie, nu mi se pare că ar îmbunătății colaborarea,
însă personal mă ajută să îmi dau seama de task-urile mai importante, să adaug
notițe pe acestea prin comentarii și să înregistrez numărul de ore cât am lucrat la ele.

Elementele de colaborare sunt de asemenea destul de bune, pot să văd cu ușurință
la ce lucrează alte persoane. Însă Jira încearcă să rezolve problemele administrării
de proiecte, astfel are multe funcții dedicate pentru asta, de exemplu Sprint-uri sau
board-urile Kanban. Din acest motiv, singurul lucru care iesă în evidență atunci
când mă gândesc la Jira sunt sistemul de log work și dashboard-urile customizabile.

\subsection{Google Keep}

O aplicație de mementouri și notițe sub diferite forme. Aceasta este foarte simplistă,
la fel cum sunt majoritatea aplicațiilor de acest tip. Un lucru important la aceasta
este faptul că este accesibilă pe majoritatea dispozitivelor întrucât se prezintă
ca aplicație web dar și nativă de mobil, iar atunci când vorbim de o aplicație de notițe
am vrea să putem să o accesăm oriunde.

Aceasta se concentrează strict pe notițe simple, din acest motiv este limitată. În interiorul
notițelor este dificil să navigăm iar gruparea notițelor se poate face prin adăugarea
unor etichete simple, ceea ce funcționează însă nu este ideal într-o aplicație complexă,
iar posibilitatea de a împărții un task în mai multe sub-task-uri din Jira este lipsită.

\subsection{Todoist}

O aplicație aparent simplă de organizare însă care încearcă să se plaseze ca fiind folositoare
și pentru echipe. Din acest motiv, task-urile pot fi împărțite pe proiecte și ulterior partajate.
Aplicația fie are un design de complex sau nu am petrecut suficient de mult timp cu ea,
însă este dificil de navigat în interiorul ei, iar pentru o aplicație de acest tip este
esențial un design minimalist.

Aceasta are mai multe funcționalități ce o apropie de Jira, precum posibilitatea de a avea
subtask-uri la task-uri, asignarea de priorități, remindere, gruparea pe proiecte
și etichete, existența unui jurnal de activitate ce descrie activitățile pe care le-am făcut
în cadrul aplicației și posiblitatea de distribuire. Prezintă un calendar în interior,
ceea ce e esențial pentru a vedea modul de distribuție a task-urilor.

\section{Cerințe}

Proiectele software urmează o analiză detaliată a specificațiilor pentru a ne
asigura că soluția pe care o implementăm rezolvă problema, însă prefer ca aceasta
să fie cât mai simplă.

Astfel, aplicația trebuie să poată să îndeplinească următoarele cerințe:

\begin{itemize}
\item Aplicația va fi disponibilă pe mai multe dispozitive și va permite sincronizarea
între dispozitive
\item Aplicația va funcționa integral fără conexiune la Internet și implicit fără niciun server
\item Aplicația va avea o componentă creată cu o arhitectură bazată pe microservicii
    \item Ca și utilizator pot să am posibilititatea de a crea un task pe care vreau să îl completez
\item Ca și utilizator pot să adaug un subtask la un task creat
    \item Ca și utilizator pot să adaug timp petrecut asupra unui task în orice moment
    \item Aplicația va fi minimalistă și adaptabilă, astfel utilizatorul poate să
își creeze un mod propriu în care aceasta să fie
\item Aplicația permite să creăm o planificare a zilei curente realizat din task-urile pe care le-am introdus
    \item Aplicația va avea o funcționalitate de planificare automată a zilei inteligentă, în funcție de tipul task-ului însă
care nu e necesar să respecte condiția legată de conexiunea de Internet
\end{itemize}

\section{Arhitectura}

Centrul aplicației va fi cluster-ul de Kubernetes. Pentru a ușura dezvoltarea vom folosii
un cluster administrat de un provider de cloud, în cazul meu Digital Ocean Managed Kubernetes.
Astfel pot să administrez un cluster cu ușurință și să modific numărul de noduri în orice moment
fără a fi nevoit să fac vreo modificare manuală în afara centrului de control.

Pentru a lansa microserviciile în Kubernetes va trebui ca acestea să fie stocate
în interiorul unui registru de containere, astfel o să creez unul tot pe Digital Ocean, având
beneficiul că nu va trebui să adminstrez modul în care administrez cheile de acces către acestea,
cluster-ul primind implicit acces către el prin diferite funcții oferite de provider.

Întrucât Kubernetes nu excelează în hostarea de baze de date voi încerca să deleg cât mai multe
funcționalități ce pot fi hostate în exterior cu putință. Astfel, voi folosii o bază de date MongoDB hostată
pe Mongo Atlas și pentru metrici și observabilitate voi folosii serviciile oferite de Grafana Cloud.

În interiorul clusetului de Kubernetes voi rula toate microserviciile prin facilitățile pe care
acesta ni le oferă alături de alte obiecte de care am avea nevoie, precum un message broker sau colector de log-uri.

Aplicația va fi dezvoltată în Flutter pentru a avea un singur code base ce duce la
crearea aplicației pe mai multe dispozitive.

\section{Implementare}

\subsection{Interfață}

În momentul în care am pornit implementarea, am crezut că cea mai mare dificultate
pe care o să o am în proiect va fi administrarea clusterului de Kubernetes, provenit din experiențele
precedente, astfel am decis că va fi o idee bună în a începe cu interfața.

Interfața principală va fi aplicația creată în Flutter. O dificultate pe care am avut-o
a fost faptul că deși eu știam în mare măsură ce vreau ca aplicația să facă, eu nu am definit
modul cum aș vrea ca aceasta să arate. Întrucât pe parcursul programului de licență nu am avut
proiecte care să se axeze pe aspectul și design-ul unei aplicații, am ajuns la un obstacol
în momentul în care am început implementarea. În mod inițial, scopul aplicației a fost să
creez ceva ce aș putea să folosesc în continuare (de aici vine cerința de a nu avea
dependințe pe niște servere, întrucât nu aș putea să le hostez mereu), însă nefiind satisfăcut
de modul cum afișez interfața a trebuit să renunț la obiectivul de a face ceva util, însă totuși
va putea fi utilizată într-o anumită măsură, deși nu ar exista neapărat un motiv pentru asta.

Va exista și o interfață ce are cu rol în autentificare și administrare a contului. O abordare generală
pe care o putem observa în aplicațiile pe care le folosim este că ele apelează același serviciu
prin care se autentifică, de exemplu când vrem să ne conectăm la un produs Google, mereu
vom fi redirecționați către accounts.google.com unde vom face autentificarea. În același mod,
atunci când vrem să facem o modificare asupra unui cont (de exemplu schimat de parolă sau a unui mail)
nu o vom facem din aplicația propriu-zisă ci tot de pe acel site de management.

Am vrut să replic acest comportament, deci va exista o interfață web minimalistă, implementată în Svelte
ce va rula în interiorul clusterului prin care vom interactiona cu microserviciul de autentificare.

\subsection{Autentificare}

Autentificarea în aplicație se va face prin intermediul microserviciului de autentificare. Acesta este
conectat la o bază de date MongoDB hostată pe Mongo Atlas unde va reține datele de care are nevoie.

Vom putea face autentificarea prin intermediul interfeței web, care atunci când se conectează cu succes
va înapoia un JSON Web Token ce conține datele necesare pentru celelalte microservicii, iar atunci când utilizatorul
încearcă să intre pe o cale protejată va trebui să prezinte acest token.

Acesta a fost creat în Go folosind framework-ul Fiber, ce încearcă să replice funcționalitatea
framework-ului Express pentru Node.JS.

Modul de acces cu acest microserviciu sunt:

\begin{itemize}
      \item POST /api/signup, pentru a crea un utilizator nou folosind mail-ul si parola
      \item POST /api/login, pentru a ne autentifica, autentificarea este folosita cu JSON Web Token.
            Fiecare cale ce necesita autentificare trebuie sa primeasca un token
            ce poate fi decodat pentru a identifica utilizatorul. Secretul pentru
            decodificare este transmis prin secrete Kubernetes catre toate microserviciile
            ce au nevoie de el.
      \item PATCH /api/updateSelf, ruta protejată folosită pentru a actualiza date non senzitive precum poza de profil.
            în momentul în care această rută este apelată cu succes, se va trimite prin Rabbit MQ evenimentul că acel utilizator și-a schimbat datele.
      \item GET /api/self, ruta protejată folosită pentru ca un utilizator să afle date despre el
      \item GET /api/users/:userid/, rută folosită pentru a afla
            datele despre un utilizator cu un anumit ID, aceasta este folosită de alte microservicii,
            astfel aceasta rută ar trebui să fie accesibilă doar din interiorul clusterului.
      \item PATCH /api/updatePassword, rută protejată folosită pentru schimbarea parolei,
            atunci când parola este schimbată acest detaliu se reflectă în baza de date și JSON Web Token-urile create anterior
            vor fi invalidate întrucât se verifică timpul emiterii token-ului și data la care a fost schimbată parola.
            Acest lucru va avea efect doar asupra acestui microserviciu, alte microservicii nu vor cunoaște același detaliu,
            astfel o metodă de invalidare ar fi prin trimiterea acestea informații pe Event Bus și fiecare microserviciu va reține un număr de tokeni invalizi.
\end{itemize}

\subsection{Sincronizarea}

Pentru a îndeplinii cerința de sincronizare pe mai multe dispozitive, va trebui creat un microserviciu
ce are ca scop memorarea datelor utilizatorului. În crearea acestei funcționalități
am întâmpinat dificultăți în crearea unei sincronizări adevărate întrucât logica
era deja implementată în interfață atunci când am ajuns la acest pas.

Astfel, pe dispozitivele mobile, datele sunt menținute într-un fișier prin intermediul SQLite3, astfel,
serviciul de sincronizare reprezintă doar o modalitate de a urca acest fișier cu date. Acest lucru
nu va duce la performanțe bune, dar va ușura dezvoltarea.

Microserviciu scris în Node.JS prin TypeScript folosind Express.
Modul cum interacționăm cu acest microserviciu:

\begin{itemize}
      \item PATCH /storage/sync/, rută protejată prin care putem să încărcăm o bază de date sub formă de fișier,
            în momentul accesării se va trimite evenimentul pe Event Bus că utilizatorul și-a modificat baza de date
      \item GET /storage/icalself, rută protejtă prin care un utilizator află link-ul către agenda lui încărcată sub formă de iCal
      \item GET /storage/ical/:id/, rută publică prin care se obține fișierul iCal pentru un anumit utilizator. Acesta poate fi folosit pentru a sincorniza aplicația cu alte aplicații de calendar.
      \item GET /storage/db/, rută protejată prin care un utilizator își poate descărca baza de date încărcate
      \item GET /storage/getdb/:id/, rută publică prin care alte microservicii pot face rost de baza de date a unui utilizator, această rută nu ar trebui să fie publică și să fie accesibilă doar în interiorul clusterului.
\end{itemize}

\subsection{Timeline}

Nu reprezenta o cerință, însă pentru a ilustra comunicarea prin evenimente am creat un Timeline, asemănător cu cel
de pe Jira, în care atunci când cineva își actualizează un task, va fi vizibil pentru toți într-un Activity Log.

Am creat un microserviciu care va memora persoanele care sunt prietene cu celelalte persoane (sau un sistem de urmărie).
În cadrul acestuia, se vor memora ID-urile persoanelor comune. La un anumit interval acest microserviciu își va actualiza datelele,
astfel va face o cerere către microserviciul de sincronizare pentru a salva intern anumite date despre task-urile la care s-a lucrat.

Microserviciu creat în Go folosind Fiber. Acesta reacționează la evenimentele de actualizare
de date ale utilizatorului și își va actualiza starea internă.

Rute disponibile:
\begin{itemize}
      \item GET /timeline/friends, rută protejată în care putem să preluăm persoanele care ne vor apărea în timeline
      \item GET /timeline/friends/:id/, rută publică în care putem să preluăm persoanele care vor apărea în timeline-ul unui anumit utilizator,
            această rută ar trebui să fie accesibilă doar în interiorul clusterului
      \item GET /timeline/, rută protejată în care putem să luăm timeline-ul propriu, timeline-ul
            corespunde din activitățile tuturor persoanelor pe care le urmărim.
      \item GET /timeline/of/:id/, rută publică în care putem să luăm timeline-ului unui anumit utilizator
      \item POST /timeline/friends, rută protejată în care adăugăm o persoană în lista de persoane urmărite
\end{itemize}

\subsection{Profile}

Pentru a complementa microserviciul de Timeline, am creat un microserviciu ce generază profilul unor persoane.
Pentru a genera profilul cuiva, am avea nevoie de informații ce provin din microserviciul de autentificare pentru
datele acestuia precum nume sau adresa de email dar și din cel de Timeline pentru a vedea toate persoanele
cu care e prieten. Astfel, se fac mai multe cereri decât am avea nevoie.

Implementarea mea ilustrează arhitectura Backend For Frontend în care am construit un microserviciu
care face cererile necesare și împachetează rezultatele pentru a efectua cererile.

Modul în care interacționăm cu acest microserviciu:
\begin{itemize}
      \item GET /profile/:id/, preluăm profilul utilizatorului,
            acesta include datele sale, elementele sale din timeline și persoanele urmărite. Aceste date sunt obținute interogând celelalte microservicii, astfel utilizatorul face un singur request ce declanșează mai multe cereri.
\end{itemize}

\subsection{Comunicarea intre microservicii}

Microserviciile vor fi lansate în cadrul Kubernetes, astfel acestea vor comunica 1 la 1 prin intermediul
DNS-ului ce funcționează în interiorul cluster-ului pe care Kubernetes îl administrează.

Va exista un message broker Rabbit MQ ce are ca scop distribuirea evenimentelor precum atunci când
un utilizator schimbă informații despre contul acestuia sau atunci când își sincronizează baza de date.
Toate microserviciile se vor înregistra pe acest message broker atunci când sunt inițializate.

În acest mod, celelalte microservicii își pot actualiza datele pe care le memorează,
de exemplu cel de timeline va face o cerere din nou pentru a actualiza datele utilizatorului ce își actualizează baza de date.

\subsection{Accesul în Cluster}

Folosind un NGINX Ingress Controller vom controla accesul. Acesta are rolul de ruta traficul ce intra în el,
în funcție de path-ul cu care este accesat. Pentru a face asta creăm un domeniu pe Name.com, creăm un certificat
folosind Let's Encrpyt pe care îl încărcăm sub formă de secret în cluster. Întrucât rulăm cluster-ul pe Digital Ocean,
creăm un Load Balancer prin care îl accesăm. Domeniul înregistrat va redirecționa către acest IP ce va face accesul în cluster.

Ingress-ul poate fi configurat cu diferite rute, sub diferite forme incluzând prefix, astfel
dacă acesta este accesat acesta se va uita după prefixul cererii și îl va redirecționa către serviciul corespunzător.
Astfel, dacă sunt accesate anumite resurse nepermise le putem interzice, întrucât Ingress-ul este folosit strict pentru comunicarea cu
exteriorul.

\section{Automatizări}

Codul va fi versionat prin intermediul GitHub, astfel folosind GitHub Actions vom putea crea
automatizări în funcție de lucrurile ce se întâmplă asupra codului nostru.

Astfel, putem să:

\begin{itemize}
    \item Atunci când publicăm un tag nou, se va crea un build al aplicației de Flutter pentru Android
          (reprezintă cea mai simplu de implementat) ce va fi stocat în GitHub Releases
    \item Atunci când publicăm modificări asupra codului unui microserviciu, iar acesta intră în branch-ul principal,
          se va declanșa un build al imaginii care va fi urcat automat in repository-ul de pe Digital Ocean
\end{itemize}

Ca o completare ale acestor automatizări putem să creăm alte diferite automatizări ce ne-ar ajuta precu atunci când modificăm
fișierele de Kubernetes, acestea să fie aplicate automat pe cluster, dar și
să implementăm teste și alte verificări prin care să ruleze în perioada înainte de publicare în Container Registry.


\section{Analiza post implementare}

Începând dezvoltarea cu frontend-ul propriu zis (aplicația în Flutter), prima dificultate a fost
incapacitatea de a traduce cerințele aplicației în perspectivă de utilizare. Desigur, puteam să
ader la o interfață simpla (fără decorații, spațieri, etc.) însă principalul obiectiv personal
a fost să creez o aplicație pe care aș putea să o folosesc iar problema UI-ului a
reprezentat problema care a dus la necompletarea acestui obiectiv întrucât nu am putut
să îmi imaginez o interfață funcțională și atractivă.

O a doua dificultate a fost faptul că înainte de a implementa interfața nu am creat o arhitectură și
o structură pentru backend, astfel pentru a manipula informații în UI, eu am creat clase
ce nu au putut fi traduse în backend, ce a dus la limitarea modului de implementare a backend-ului.

Implementarea backend-ului a presupus crearea unui cluster local de Kubernetes,
întrucât nu am distribuit acest mediu eficient, orice modificare a însemnat un delay de 5-10 minute
(modificări în cod, push pe GitHub, build de imagine, urcat în repository, actualizare la fișiere
de Kubernetes sau kill la Pod pentru recreere).

Platforma pe care am dezvoltat în principal a fost un telefon cu Android, pentru a accesa clusterul local
care rula pe laptop a trebuit să le conectez la același WiFi, să utilizez o aplicație asemănătoare cu un VPN (
avea rolul de a monitoriza traficul) în care am făcut setări pentru a direcționa adresa IP locală a laptopului
către domeniul setat în Ingress Controller. Această dificultate a făcut dificilă testarea pe alte platforme (precum un dispozitiv iOS).

O variantă mai bună pentru încărcarea fișierului întreg de SQL pe discul microserviciului de stocare
(fapt ce limitează scalarea acestuia, putând fii doar unul) ar fi utilizarea de Blob Storage, astfel
încărcam baza de date acolo, dar microserviciul ar fi fucnționat ca un intermediar ce putea să cacheze aceste baze de date.

Regenerarea certificatului necesar pentru TLS în Kubernetes putea fi făcută automat folosind un CronJob, astfel
aceasta a fost necesară să fie făcută manual.


\chapter{Concluzii}

O arhitectură bazată pe microservicii reprezintă o modernizare a modului de a dezvolta aplicații.
Cuplat cu platformele de infrastructură în cloud ce ne permit crearea de servere
în orice moment într-un timp foarte scurt și diferitele tipuri de orchestrare
putem crea aplicații scalabile și reziliente.

Microserviciile sunt o arhitectură de dezvoltare a aplicațiilor în care aplicația este
împărțită în componente independente ce reprezintă o parte indepententă și este responsabilă pentru o
funcționalitate specifică al procesului de afaceri. Aceste servicii comunică între
ele prin intermediul unor interfețe bine definite și pot fi dezvoltate, implementate și scalate independent.

Prin utilizarea arhitecturii bazate pe microservicii,
aplicațiile devin mai ușor de întreținut și scalat,
deoarece fiecare serviciu poate fi actualizat și scalat independent,
fără a afecta restul sistemului.
De asemenea, permite dezvoltarea agilă și
implementarea continuă, deoarece fiecare serviciu poate
fi testat și implementat separat.

Cu toate acestea, arhitectura bazată pe microservicii
introduce și provocări suplimentare în gestionarea
comunicării între servicii și în monitorizarea și
administrarea infrastructurii distribuite. Amazon Prime Video a observat
o creștere semnificativă de performană atunci când au făcut trecerea de la o arhitectură
distribuită către una monolitică \cite{Kolny_2023}. În cazul acestora, comunicațiile
dintre părțile componente au dus la consum suplimentar de resurse.

Deși nu unica metodă de orchestrare pentru microservicii, Kubernetes reprezintă o metodă complexă de orchestrare a containerelor ce oferă
un mediu scalabil și flexibil pentru implementarea aplicațiilor. Oferă diferite
funcționalități esențiale precum descoperirea în interiorul cluster-ului,
posibilitatea de a gestiona actualizări ale configurației sub formă de fișiere ce permite folosirea
acestuia împreună cu sisteme de gestiune a infrastructurii prin cod dar și
distribuirea sarcinilor în mod inteligent și rutarea traficului în interiorul clusterului.

Serviciile diversificate oferite de distribuitorii de infrastructură în cloud
accelerează dezvoltarea, astfel putem evita implementarea
unor funcționalități comune, crearea infrastructurii sau configurarea bazelor de
date la cerere. Astfel putem să ne concentrăm pe construirea de soluții și nu pe modul
cum acestea vor fi implementate.

Această lucrare abordează elemente generale legate de această arhitectură și modul ei de
orchestrare. O posibilă dezvoltare ulterioară reprezintă aprofundarea unei singure
componente prezentată și analizarea mai adâncă a diferitelor moduri în care ea impactează sistemul.
În același scop, Kubernetes reprezintă doar unul din modurile
de orchestrare ale acestei arhitecturi descris succint, având diferite dezavantaje precum complexitatea
ridicată, astfel putem să abordăm alte soluții precum Container as a Service
sau alternative mai ușoare ale acestuia precum HashiCorp Nomad.

Dezvoltarea unei aplicații pe bază de microservicii a reprezintat o dificultate ridicată,
mai ales considerând că nu a adus niciun beneficiu pentru mine. Crearea unui mediu de dezvoltare,
rularea tuturor microserviciilor pe același computer dar și stabilirea arhitecturii
au fost mai mult impedimente ale dezvoltării decât să o faciliteze. Având posibilitatea
de a folosii servicii în Cloud reprezintă unul din cel mai mare avantaj în dezvoltarea modernă,
întrucât elimină necesitatea a unei invesiții inițiale în procurarea de infrastructură
sau crearea diferitelor servicii precum monitorizare și putem să folosim obiecte deja create.


\printbibliography[heading=bibintoc]
\nocite{*}

\end{document}